\documentclass[output=paper,colorlinks,citecolor=brown]{langscibook}
\ChapterDOI{10.5281/zenodo.21237016}
\author{Jeremy Rau\orcid{}\affiliation{Harvard University}}
%\ORCIDs{}

\title{A Late Nuclear Proto-Indo-European verbal type: Intransitive thematic nasal-infix present actives}
\shorttitlerunninghead{A Late Nuclear Proto-Indo-European verbal type}

\abstract{This paper examines primary verbal formations in the ancient Indo-European languages that also play a role in deadjectival derivation. It builds on recent work on the historical morphology of the intransitive change of state nasal-affix thematic active presents in the “Northern Indo-European” languages – Baltic, Slavic, and Germanic – which continue an anticausative to the well-known athematic nasal-infix transitive presents. The paper argues that this formation has exact cognates in Ancient Greek, Vedic Sanskrit, and Indo-Iranian, more generally, where we find morphologically identical presents that as in the Northern Indo-European languages make semantically intransitive change of state formations to both primary verbs and root-derived property concept adjectives. The agreement across these branches supports the reconstruction of this formation for late Nuclear Indo-European, as a continuation of a deeper Proto-Indo-European type.}


\IfFileExists{../localcommands.tex}{
   \addbibresource{../localbibliography.bib}
   \usepackage{tabularx,multicol}
\usepackage{url}
\urlstyle{same}

 
\usepackage{langsci-optional}
\usepackage{langsci-lgr}
\usepackage{langsci-gb4e}
\usepackage[linguistics,edges]{forest}
\usepackage{tikz-qtree}
\usetikzlibrary{positioning}
\usetikzlibrary{patterns.meta}
\usepackage{multirow}
\usepackage{pgfplots}
\usepackage{setspace}
\usepackage{amsmath,stackengine}
% \usepackage{fontspec}
% \usepackage{tipa} % TIPA is banned
\usepackage{langsci-textipa}
\usepackage{langsci-branding}
\usepackage{longtable}
\usepackage{tabto}

\usepackage{enumitem}

%\usepackage{lscape} %Querformat chapter 4

   % \newcommand*{\orcid}{}


\makeatletter
\let\thetitle\@title
\let\theauthor\@author
\makeatother

\newcommand{\togglepaper}[1][0]{
   \bibliography{../localbibliography}
   \papernote{\scriptsize\normalfont
     \theauthor.
     \titleTemp.
     To appear in:
    Laura Grestenberger, Viktoria Reiter \& Melanie Malzahn (eds.).
    \emph{Deadjectival verb formation in Indo-European and beyond}.
     Berlin: Language Science Press. [preliminary page numbering]
   }
   \pagenumbering{roman}
   \setcounter{chapter}{#1}
   \addtocounter{chapter}{-1}
}

\newbool{bookcompile}
\booltrue{bookcompile}
\newcommand{\bookorchapter}[2]{\ifbool{bookcompile}{#1}{#2}}

%Sonderzeichen
%idg palatales k
\newcommand{\kinvbreve}{k\raisebox{0.6ex}{\hspace{-0.05em}\char"0311}}

\newcommand{\jac}{%
  \textit{\char"0237}% italic dotless j (ȷ)
  \hspace{-0.01em}% fine-tune horizontal spacing
  \raisebox{0.01ex}{\char"0301}% raised combining acute
}

%Baltic circumflex l
\newcommand\ltilded{%
\hspace{-0.2em}%
  \stackon[-1.5pt]{l}{\smash{\kern0.2em\~{}}}%
  \hspace{-0.2em}%
}

%Slav. wedge+ double acute
\newcommand{\HighH}[1]{%
    \stackon[-1.2ex]{#1}{\kern0.1em\H{}\kern-0.1em}%
}

%Lith dot tilde VICKY diese b.-sl. Sonderzeichen BRINGEN MICH UM!!
\newcommand{\tildedot}[1]{%
    \ensurestackMath{%
        \stackon[-0.27ex]{% Tilde layer 
            \stackon[-0.15ex]{% Dot layer 
                #1%
            }{\hspace{0.1em}\cdot\hspace{-0.13em}}%
        }{\hspace{0.1em}\text{\~{}}\hspace{-0.13em}%
    }%
}}

%Lith dot acute
\newcommand{\dotacute}[1]{%
    \ensurestackMath{%
        \stackon[-1.1ex]{% acute layer 
            \stackon[-0.15ex]{% Dot layer 
                #1%
            }{\hspace{0.1em}\cdot\hspace{-0.12em}}%
        }{\hspace{0.1em}\text{\'{}}\hspace{-0.12em}%
    }%
}}



% \setlength{\emergencystretch}{2pt} 

 \pgfplotsset{compat=1.18}


\defcitealias{AbB_1}{AbB 1}
\defcitealias{AbB_7}{AbB 7}
\defcitealias{ABL_1}{ABL 1}
\defcitealias{Vries1977}{AEW}
\defcitealias{ALEW}{ALEW}
\defcitealias{AnGr}{AnGr}
\defcitealias{ARM_1}{ARM 1}
\defcitealias{BT}{BT}
\defcitealias{black_concise_2000}{CDA}
\defcitealias{CHD}{CHD}
\defcitealias{CT28}{CT 28}
\defcitealias{CT34}{CT 34}
\defcitealias{CT39}{CT 39}
\defcitealias{Cleasby1874}{CV}
\defcitealias{Beekes2010}{EDG}
\defcitealias{eDiAna}{eDiAna}
\defcitealias{eDIL}{eDIL}
\defcitealias{EDL}{EDL}
\defcitealias{Kroonen2013}{EDPG}
\defcitealias{ESJS}{ESJS}
\defcitealias{EVP}{EVP}
\defcitealias{EWA1988}{EWA}
\defcitealias{EWAI}{EWAia I}
\defcitealias{EWAII}{EWAia II}
\defcitealias{EWGP}{EWGP}
\defcitealias{GDLI1961}{GDLI}
\defcitealias{GPC}{GPC}
\defcitealias{GRADIT1999}{GRADIT}
\defcitealias{HW}{HW}
\defcitealias{IEW}{IEW}
\defcitealias{KAH_2}{KAH 2}
\defcitealias{KAR_1}{KAR 1}
\defcitealias{KAR_2}{KAR 2}
\defcitealias{KEWA}{KEWA}
\defcitealias{Labat_TDP}{Labat TDP}
\defcitealias{lambert_BWL}{Lambert BWL}
\defcitealias{LEIA}{LEIA}
\defcitealias{Rix2001}{LIV\textsuperscript{2}}
\defcitealias{LIVaddenda}{LIV\textsuperscript{2} {A}ddenda}
\defcitealias{LKG1971}{LKG}
\defcitealias{LKZ}{LKŽ}
\defcitealias{LP}{LP}
\defcitealias{MhdGr}{MhdGr}
\defcitealias{MLLVG1959}{MLLVG}
\defcitealias{NIL}{NIL}
\defcitealias{OgnS}{OgnS}
\defcitealias{ONP}{ONP}
\defcitealias{RIMA_2}{RIMA 2}
\defcitealias{TIM_2}{TIM 2}
\defcitealias{WAP}{WAP}
\defcitealias{YOS_2}{YOS 2}






\newindex{wan}{wdx}{wnd}{Form index}
\newcommand{\iw}[2]{\index[wan]{#1!{#2}\xspace}}
\newcommand{\iwi}[3][]{\index[wan]{#2!{#3}\xspace#1}{#3}}
\newcommand{\iwit}[3][]{\rule{0pt}{0pt}#1\index[wan]{#2!{#3}}{#3}}


\newindex{van}{vdx}{vnd}{Form index}
\newindex{ban}{bdx}{bnd}{Book index}

% \newcommand{\iw}[3][]{\index[wan]{#2!\textit{#3}}}


% \providecommand{\iwi}[3][]{\iw{#3}\textit{#3}}


\newcommand{\indexnote}[1]{{\footnotesize\sffamily\raggedright\normalfont\color{gray}#1}}

   %% hyphenation points for line breaks
%% Normally, automatic hyphenation in LaTeX is very good
%% If a word is mis-hyphenated, add it to this file
%%
%% add information to TeX file before \begin{document} with:
%% %% hyphenation points for line breaks
%% Normally, automatic hyphenation in LaTeX is very good
%% If a word is mis-hyphenated, add it to this file
%%
%% add information to TeX file before \begin{document} with:
%% %% hyphenation points for line breaks
%% Normally, automatic hyphenation in LaTeX is very good
%% If a word is mis-hyphenated, add it to this file
%%
%% add information to TeX file before \begin{document} with:
%% \include{localhyphenation}
\hyphenation{
    par-a-digm non-de-adj-ect-iv-al -pre-sent -pre-sents for-ma-tion-al in-do-ger-ma-ni-sche in-do-ger-ma-ni-schen Rei-chert Rei-ter dia-chrony ety-mo-lo-gi-sches Ost-ger-ma-ni-schen ger-ma-ni-schen alt-in-di-schen mor-pho-lo-gi-scher Alt-indo-ari-sch-en lexico-graphy Schrij-ver
}
% aktuell in Bibliographie: stud-ies, histor-ical Soci-età
\hyphenation{
    par-a-digm non-de-adj-ect-iv-al -pre-sent -pre-sents for-ma-tion-al in-do-ger-ma-ni-sche in-do-ger-ma-ni-schen Rei-chert Rei-ter dia-chrony ety-mo-lo-gi-sches Ost-ger-ma-ni-schen ger-ma-ni-schen alt-in-di-schen mor-pho-lo-gi-scher Alt-indo-ari-sch-en lexico-graphy Schrij-ver
}
% aktuell in Bibliographie: stud-ies, histor-ical Soci-età
\hyphenation{
    par-a-digm non-de-adj-ect-iv-al -pre-sent -pre-sents for-ma-tion-al in-do-ger-ma-ni-sche in-do-ger-ma-ni-schen Rei-chert Rei-ter dia-chrony ety-mo-lo-gi-sches Ost-ger-ma-ni-schen ger-ma-ni-schen alt-in-di-schen mor-pho-lo-gi-scher Alt-indo-ari-sch-en lexico-graphy Schrij-ver
}
% aktuell in Bibliographie: stud-ies, histor-ical Soci-età
   \boolfalse{bookcompile}
   \togglepaper[9]%%chapternumber
}{}
\setlength{\emergencystretch}{2pt} 
\begin{document}
\SetupAffiliations{mark style=none}
\maketitle

\section{Introduction} 

 


One of the more remarkable features of Indo-European (IE) morphology is that deadjectival verbal\is{deadjectival!verb} formations are partially handled with dedicated denominative morphology and partially with the root-derived\is{derivation!deradical} -- ``primary'' -- morphology that otherwise holds across Proto-Indo-European (PIE) verbal classes. While the deadjectival denominative formations\is{denominative!verb} have been thoroughly studied and are well understood, the insight that primary formations\is{derivation!deradical} also played a role in deadjectival derivation\is{derivation!deadjectival}  -- at least in the sphere of property concept adjectives -- is a recent discovery whose full parameters have yet to be elucidated.\footnote{See \citealt{Rau2009} with lit., \citeyear{Rau2013}, \citeyear{Rau2016}, \citeyear{Rau2017}, and further \citealt{Nussbaum2022}.}  

This paper aims to further advance this discussion and takes its start from recent work on the historical morphology of the \isi{intransitive} change of state\is{change of state (CoS)} nasal-affix thematic active presents \is{nasal present} in the “\isi{Northern Indo-European}” languages (Baltic, Slavic, and Germanic), whose origins have been suggested to continue a morphologically peculiar \isi{anticausative} to the well-known athematic nasal-infix \isi{transitive} presents.\is{nasal present} The paper argues that this type has exact cognates in Ancient Greek and Vedic Sanskrit (and Indo-Iranian, more generally), where we find morphologically identical presents that as in the \isi{Northern Indo-European} languages make semantically \isi{intransitive} change of state\is{change of state (CoS)} formations to both primary verbs\is{derivation!deradical} and subsets of root-derived\is{derivation!deradical} property concept\is{property concept (PC)} adjectives. Based on the agreement of these branches, it is possible to reconstruct the formation for late Nuclear Indo-European\footnote{“Nuclear Indo-European” is used to refer to the subgroup of Indo-European languages that shared innovations after the departure of Anatolian from the PIE speech community. I use “late Nuclear PIE” to refer to the subgroup of Nuclear PIE languages that remain after the departure of Tocharian and then later Proto-Italo-Celtic.} and to project its origins deep into PIE.

\section{Northern Indo-European}

The starting point for this investigation is the magisterial 2007 Harvard PhD thesis of Yaroslav Gorbachov (\citealt{Gorbachov2007}; see further  \citealt{VillanuevaSvensson2006}, \citeyear{VillanuevaSvensson2011}) on the \isi{inchoative} nasal-affix presents \is{nasal present} in Baltic, Slavic, and Germanic, the \isi{Northern Indo-European} languages. This nasal-affix \is{nasal present} verbal type is well attested in all three \isi{Northern Indo-European} branches, where it generally serves to make 1) \isi{intransitive} \is{change of state (CoS)} change of state/location formations -- often categorized as “\isi{inchoative}s” -- to \isi{intransitive} \isi{stative}s, less frequently \isi{activity} verbs, 2) \isi{anticausative}s to \isi{causative} formations, and 3) denominatives\is{denominative!verb} to property concept\is{property concept (PC)} adjectives and, more rarely, substantives. For a survey of forms across all representative semantic categories, cf.\ the examples in (\ref{ex:inchoa}).

\begin{exe}
\ex Inchoative nasal-affix presents in Slavic, Baltic, and Germanic \label{ex:inchoa}
\begin{xlist}
    \ex \textbf{Old Church Slavonic (OCS)}
    \begin{itemize}
        \item \textit{sędǫ, sesti, sědŭ}\iw{OCS}{\textit{sesti, sędǫ, sědŭ}} ‘sit down’
        \item \textit{vŭzlęgǫ, -lešti, -legŭ}\iw{OCS}{\textit{vŭzlešti, -lęgǫ, -legŭ}} ‘lie down’
        \item \textit{stanǫ, stati, staxŭ}\iw{OCS}{\textit{stati, stanǫ, staxŭ}} ‘stand up’
        \item \textit{pri/vŭzniknǫ, -niknǫti, -nikŭ}\iw{OCS}{\textit{pri/vŭzniknǫti, -niknǫ, -nikŭ}} ‘spring up’
        \item \textit{promŭknǫ, -mŭknǫti sę, -mŭkŭ}\iw{OCS}{\textit{promŭknǫti sę, -mŭknǫ, -mŭkŭ}} ‘go out, spread’
        \item \textit{prilĭ(p)nǫ, -lĭ(p)nǫti, -lĭpŭ}\iw{OCS}{\textit{prilĭ(p)nǫti, -lĭ(p)nǫ, -lĭpŭ}} ‘cleave, cling to’
        \item \textit{uglĭbnǫ, -glĭbnǫti, -glĭbŭ}\iw{OCS}{\textit{uglĭbnǫti, -glĭbnǫ, -glĭbŭ}} ‘get stuck’
        \item \textit{izběgnǫ, -běgnǫti, -běgŭ}\iw{OCS}{\textit{izběgnǫti, -běgnǫ, -běgŭ}} ‘flee, avoid’
        \item \textit{postignǫ, -stignǫti, -stigŭ}\iw{OCS}{\textit{postignǫti, -stignǫ, -stigŭ}} ‘come upon, reach’
        \item \textit{vŭzbŭ(d)nǫ, -bŭ(d)nǫti, -bŭdŭ}\iw{OCS}{\textit{vŭzbŭ(d)nǫti, -bŭ(d)nǫ, -bŭdŭ}} ‘wake up’
        \item \textit{usŭ(p)nǫ, -sŭ(p)nǫti, -sŭpŭ}\iw{OCS}{\textit{usŭ(p)nǫti, -sŭ(p)nǫ, -sŭpŭ}} ‘fall asleep’
        \item \textit{navyknǫ, navyknǫti, navykŭ}\iw{OCS}{\textit{navyknǫti, navyknǫ, navykŭ}} ‘get used to’
        \item \textit{oslĭpnǫ, -slĭpnǫti, -slĭpŭ}\iw{OCS}{\textit{oslĭpnǫti, -slĭpnǫ, -slĭpŭ}} ‘go blind’ (: \iwi{OCS}{\textit{slěpŭ}})
        \item \textit{ugasnǫ, -gasnǫti, -gasŭ}\iw{OCS}{\textit{ugasnǫti, -gasnǫ, -gasŭ}} ‘go out, be extinguished’
        \item \textit{prisvędnǫ, -svędnǫti, -svędŭ}\iw{OCS}{\textit{prisvędnǫti, -svędnǫ, -svędŭ}} ‘wilt’
        \item \textit{omrĭknǫ, -mrĭknǫti, -mrĭkŭ}\iw{OCS}{\textit{omrĭknǫti, -mrĭknǫ -mrĭkŭ}} ‘grow dark’ (: \iwi{OCS}{\textit{mĭrkŭ}*})
        \item ORuss.\ \iwi{ORuss}{\textit{svĭ(t)nuti}} ‘grow light’
        \item \textit{sŭ/pomrŭznǫ, mrŭznǫti, -mrŭzŭ}\iw{OCS}{\textit{su/pomrŭznǫti, -mrŭznǫ, -mrŭzŭ}} ‘freeze, congeal’ (: \iwi{OCS}{\textit{mrazŭ}})
        \item \textit{sŭxnǫ, sŭxnǫti, sŭxŭ}\iw{OCS}{\textit{sŭxnǫti, sŭxnǫ, sŭxŭ}} ‘dry up’ (: \iwi{OCS}{\textit{suxŭ}})
    \end{itemize}
    \ex \textbf{Baltic}
    \largerpage
    \begin{itemize}
        \item OPr.\ \iwi{OPr}{\textit{sīnda}-} ‘sit down’
        \item OPr.\ \textit{polīnka, polāikt}\iw{OPr}{\textit{polāikt, polīnka}} ‘stays’ (: Lith.\ dial. -\textit{liñka} ‘remains’)
        \item Lith.\ \textit{tam̃pa, tàpti, tàpo}\iw{Lith}{\textit{tàpti, tam̃pa, tàpo}} ‘become’
        \item Lith.\ \textit{knim̃ba, knìbti, knìbo}\iw{Lith}{\textit{knìbti, knim̃ba, knìbo}} ‘sink, fall’
        \item Lith.\ \textit{kañka, kàkti, kàko}\iw{Lith}{\textit{kàkti, kañka, kàko}} ‘start out, go’
        \item Lith.\ \textit{muñka, mùkti, mùko}\iw{Lith}{\textit{mùkti, muñka, mùko}} ‘break free, escape’
        \item Lith.\ \textit{lim̃pa, lìpti, lìpo}\iw{Lith}{\textit{lìpti}, \textit{lim̃pa}, \textit{lìpo}} ‘cling, stick to, climb’
        \item Lith.\ \textit{pliñta, plìsti, plìto}\iw{Lith}{\textit{plìsti, pliñta, plìto}} ‘spread out’
        \item Lith.\ \textit{skiñda, skìsti, skìdo}\iw{Lith}{\textit{skìsti, skiñda, skìdo}} ‘fall apart, split asunder’
        \item Lith.\ \textit{atbuñda, atbùsti, atbùdo}\iw{Lith}{\textit{atbùsti, atbuñda, atbùdo}} ‘wake up’
        \item Lith.\ \textit{miñga, mìgti, mìgo}\iw{Lith}{\textit{mìgti, miñga, mìgo}} ‘fall asleep’
        \item Lith.\ \textit{suprañta, supràsti, supràto}\iw{Lith}{\textit{supràsti, suprañta, supràto}} ‘understand’
        \item Lith.\ \textit{patiñka, patìkti, patìko}\iw{Lith}{\textit{patìkti, patiñka, patìko}} ‘suits, fits’
        \item Lith.\ \textit{šviñta, švìsti, švìto}\iw{Lith}{\textit{švìsti, šviñta, švìto}, 1sg. prs. \textit{švintù}} ‘dawn’
        \item Lith.\ \textit{sniñga, snìgti, snìgo}\iw{Lith}{\textit{snìgti, sniñga, snìgo}} ‘snow’
        \item Lith.\ \textit{tuñka, tùkti, tùko}\iw{Lith}{\textit{tùkti, tuñka, tùko}} ‘grow fat’ (: \iwi{Lith}{\textit{taukaĩ}})
        \item Lith.\ \textit{pųva, pùti, pùvo}\iw{Lith}{\textit{pùti, pųva, pùvo}} ‘rot, decay’
        \item Lith.\ \textit{sųsa, sùsti, sùso}\iw{Lith}{\textit{sùsti, sųsa, sùso}} ‘dry up’ (: \iwi{Lith}{\textit{saũsas}})
        \item Lith.\ \textit{dum̃ba, dùbti, dùbo}\iw{Lith}{\textit{dùbti, dum̃ba, dùbo}} ‘grow hollow’ (: \iwi{Lith}{\textit{dubùs}})
    \end{itemize}
    \ex \textbf{Germanic}
    \begin{itemize}
        \item Go.\ \iwi{Goth}{\textit{standan}} ‘stand', ON \iwi{ON}{\textit{standa}} ‘id.’
        \item Go.\ \textit{aflifnan}\iw{Goth}{\textit{aflifnan, -lifniþ, -lifnoda}} ‘remain', ON \iwi{ON}{\textit{lifna}} ‘id.’
        \item OE \iwi{OE}{\textit{climban}} ‘climb', OHG \iwi{OHG}{\textit{klimban}} ‘id.’
        \item Go.\ \iwi{Goth}{\textit{rinnan}} ‘run', ON \iwi{ON}{\textit{renna}} ‘id.’
        \item Go.\ \iwi{Goth}{\textit{aflinnan}} ‘leave off, desert', ON \iwi{ON}{\textit{linna}} ‘cease, stop’
        \item Go.\ \iwi{Goth}{\textit{wiknan}} ‘start to give in, yield', ON \iwi{ON}{\textit{vikna}} ‘id.’
        \item Go.\ \iwi{Goth}{\textit{usgutnan}} ‘spill out’
        \item Go.\ \iwi{Goth}{\textit{usbriknan}} ‘break off’
        \item Go.\ \iwi{Goth}{\textit{gawaknan}} ‘wake up', ON \iwi{ON}{\textit{vakna}} ‘id.’
        \item ON \iwi{ON}{\textit{sofna}} ‘fall asleep’
        \item Go.\ \iwi{Goth}{\textit{ga(h)nipnan}} ‘grow sad', ON \iwi{ON}{\textit{hnipna}} ‘become downcast, droop’
        \item Go.\ \iwi{Goth}{\textit{gawairþnan}} ‘become reconciled’
        \item Go.\ \iwi{Goth}{\textit{gathaursnan}} ‘get dry', ON \iwi{ON}{\textit{þorna}} ‘id.’
        \item Go.\ \iwi{Goth}{\textit{gastaurknan}} ‘grow stiff, rigid', ON \iwi{ON}{\textit{storkna}} ‘id.’ (: ON \iwi{ON}{\textit{starkr}})
        \item Go.\ \iwi{Goth}{\textit{bulgnan}} ‘swell up', ON \iwi{ON}{\textit{bolgna}} ‘id.’
        \item ON \iwi{ON}{\textit{fúna}} ‘rot’
        \item Go.\ \textit{gafullnan}\iw{Goth}{\textit{(ga)fullnan}} ‘grow full', ON \iwi{ON}{\textit{fullna}} ‘id.’ (: \iwi{Goth}{\textit{fulls}})
        \item ON \textit{hvítna}\iw{ON}{\textit{hvítna} (OWN)} ‘goes pale’ (: \iwi{ON}{\textit{hvítr}})
        \item ON \textit{róðna}\iw{ON}{\textit{roðna} (OWN)} ‘turn red’ (: \iwi{ON}{\textit{rauðr}})
    \end{itemize}
\end{xlist}
\end{exe}

This \isi{Northern Indo-European} verbal type also shows a large number of two- and three-way cognate sets among the languages -- cf.\ the forms in \tabref{tab:key:1}.

\begin{table}
\caption{Cognate inchoative nasal-affix presents in Slavic, Baltic, and Germanic}
\label{tab:key:1}
\begin{tabularx}{\textwidth}{QQQ}
\lsptoprule
{\bfseries OCS} & { \textbf{Baltic}} & { \textbf{Germanic}}\\
\midrule
{\mbox{}\textit{sędǫ, sesti, sědŭ}\iw{OCS}{\textit{sesti, sędǫ, sědŭ}} ‘sit down’} & {OPr.\ \iwi{OPr}{\textit{sīnda}-} ‘sit down’} & \\
\tablevspace
{\mbox{}\textit{prilĭ(p)nǫ, -lĭ(p)nǫti, -lĭpu} ‘cleave, cling to’}\iw{OCS}{\textit{prilĭ(p)nǫti, -lĭ(p)nǫ, -lĭpŭ}} & { Lith.\ \textit{lim̃pa, lìpti, lìpo}\iw{Lith}{\textit{lìpti}, \textit{lim̃pa}, \textit{lìpo}} ‘cling, stick to, climb’} & { Go.\ \textit{aflifnan}\iw{Goth}{\textit{aflifnan, -lifniþ, -lifnoda}} ‘remain’}\\
\tablevspace
{\mbox{}\textit{promŭknǫ, -mŭknǫti sę, \mbox{-mŭkŭ}}\iw{OCS}{\textit{promŭknǫti sę, -mŭknǫ, -mŭkŭ}} ‘go out, spread’} & { Lith.\ \textit{muñka, mùkti, mùko}\iw{Lith}{\textit{mùkti, muñka, mùko}} ‘break free’} & \\
\tablevspace
{\mbox{}\textit{uglĭbnǫ, -glĭbnǫti, -glĭbŭ}\iw{OCS}{\textit{uglĭbnǫti, -glĭbnǫ, -glĭbŭ}} ‘get stuck’} &  & { OE \iwi{OE}{\textit{climban}} ‘climb', OHG \iwi{OHG}{\textit{klimban}} ‘id.’}\\
\tablevspace
{\mbox{}\textit{vŭzbŭ(d)nǫ, -bŭ(d)nǫti, \mbox{-bŭdŭ}}\iw{OCS}{\textit{vŭzbŭ(d)nǫti, -bŭ(d)nǫ, -bŭdŭ}} ‘wake up’} & { Lith.\ \textit{atbuñda, atbùsti, atbùdo}\iw{Lith}{\textit{atbùsti, atbuñda, atbùdo}} ‘id.’} & \\
\tablevspace
{\mbox{}\textit{usŭ(p)nǫ, -sŭ(p)nǫti, -sŭpŭ}\iw{OCS}{\textit{usŭ(p)nǫti, -sŭ(p)nǫ, -sŭpŭ}} ‘fall asleep’} &  & { ON \iwi{ON}{\textit{sofna}} ‘fall asleep’}\\
\tablevspace
{ORuss.\ \iwi{ORuss}{\textit{svĭ(t)nuti}} ‘grow light’} & { Lith.\ \textit{šviñta, švìsti, švìto}\iw{Lith}{\textit{švìsti, šviñta, švìto}, 1sg. prs. \textit{švintù}} ‘dawn’} & { ON \textit{hvítna}\iw{ON}{\textit{hvítna} (OWN)} ‘goes pale’}\\
\lspbottomrule
\end{tabularx}
\end{table}

In his thesis, Gorbachov makes three important observations about this type and its prehistory. First, he notes that these verbs inflect as \textit{active} and \textit{thematic} presents \is{nasal present} with a nasal affix across all three branches -- so in Gothic and residually elsewhere in Germanic --, and that this active thematic inflectional pattern should be reconstructed for the ancestor of the type in all three languages. See \citealt{Gorbachov2007} \textit{passim}, with discussion of alternative analyses. Second, he points out that the type systematically corresponds to \isi{transitive}, athematically inflected nasal-infix presents \is{nasal present} in other IE languages, and is ultimately to be aligned with this \isi{transitive} type -- see, e.g., the forms in \tabref{tab:key:2}. 

\begin{table}[p]
\small
\caption{Correspondences between inchoative nasal-affix presents in Slavic, Baltic, and Germanic and transitive athematic nasal-infix presents in other IE languages}
\label{tab:key:2}
\begin{tabularx}{\textwidth}{QQQ>{\raggedright}p{35mm}@{}c@{}}
\lsptoprule
{ \textbf{OCS}} & { \textbf{Baltic}} & { \textbf{Germanic}} & { \textbf{Nasal-infix} \textbf{pres.}}&\\
\midrule
{\mbox{}\textit{prilĭ(p)nǫ, -lĭ(p)nǫti, -lĭpu}\iw{OCS}{\textit{prilĭ(p)nǫti, -lĭ(p)nǫ, -lĭpŭ}} ‘cleave, cling to’} & { Lith.\ \textit{lim̃pa, lìpti, lìpo}\iw{Lith}{\textit{lìpti}, \textit{lim̃pa}, \textit{lìpo}} ‘cling, stick to, climb’} & { Go.\ \textit{aflifnan}\iw{Goth}{\textit{aflifnan, -lifniþ, -lifnoda}} ‘remain’} &  Ved.\ \iwi{Sanskrit}{\textit{limpáti}} AV+, \textsuperscript{?}Gk.\ \iwi{AG}{λιπαίνω}  ‘anoint’ Att.\ Ion.&\\
\tablevspace
{\mbox{}\textit{promŭknǫ, \mbox{-mŭknǫti} sę, \mbox{mŭkŭ}}\iw{OCS}{\textit{promŭknǫti sę, -mŭknǫ, -mŭkŭ}} ‘go out, spread’} & { Lith.\ \textit{muñka, mùkti, mùko}\iw{Lith}{\textit{mùkti, muñka, mùko}} ‘break free, escape’} &  & { Ved.\ \textit{muñcáti} \iw{Sanskrit}{\textit{muñcáti/te}} `release, let go', Lat.\ \iwi{Latin}{\textit{ēmungō}}, -\textit{ere} `blow the nose'}&\\
\tablevspace
{\mbox{}\textit{vŭzbŭ(d)nǫ, \mbox{-bŭ(d)nǫti,} -bŭdŭ}\iw{OCS}{\textit{vŭzbŭ(d)nǫti, -bŭ(d)nǫ, -bŭdŭ}} ‘wake up’} & { Lith.\ \textit{atbuñda, atbùsti, atbùdo}\iw{Lith}{\textit{atbùsti, atbuñda, atbùdo}} ‘id.’} &  & { OIr.\ \iwi{OIr}{\textit{ad}$\cdot$\textit{boind}} `give notice' \mbox{(Gk.} \iwi{AG}{πυνθάνoμαι} Hom.+)}&\\
\tablevspace
& { Lith.\ \textit{skiñda, skìsti, skìdo}\iw{Lith}{\textit{skìsti, skiñda, skìdo}} ‘fall apart, split asunder’} &  & { Ved.\ \iwi{Sanskrit}{\textit{chinátti}} ‘split off, break off', Lat.\ \iwi{Latin}{\textit{scindō}} ‘id.’}&\\
\tablevspace
&  & { OHG \iwi{OHG}{\textit{scrintan}} ‘split open, burst’} & { Ved.\ \iw{Sanskrit}{\textit{kr̥ntáti/te}} \textit{kr̥ntáti} ‘cut, split', YAv.\ \iwi{Avestan}{\textit{kərəṇta}-}\textit{\textsuperscript{ti}} ‘id.’}&\\
\tablevspace
{ ORuss.\ \iwi{ORuss}{\textit{rínut’}} ‘stream, flow’} &  & { Go.\ \iwi{Goth}{\textit{rinnan}}, ON \iwi{ON}{\textit{renna}} ‘run’} & { Ved.\ \iwi{Sanskrit}{\textit{riṇ{ā́}ti}} ‘set in motion, swirl’, Gk.\ \iwi{AG}{ὀρῑ́νω}  ‘stir, move excite’ Hom.+}&\\
\lspbottomrule
\end{tabularx}
\end{table} 

Third, he argues that based on the function of the formation, its active thematic inflection, and its correlation with \isi{transitive} athematic nasal-infix verbs in the other IE languages, it is possible to set up a descriptive derivational process for the ancestor of these branches, whereby an \isi{anticausative} is made to the \isi{transitive} athematic nasal-infix present \is{nasal present} by putting the stem in the zero-grade, suffixing the accented thematic vowel, and inflecting the formation as an active, as outlined in \tabref{tab:key:3}.

\begin{table}[p]
\small
\caption{Derivational relationship between athematic nasal-infix transitives and thematic nasal-infix intransitives}
\label{tab:key:3}
\begin{tabularx}{\textwidth}{>{\raggedright\arraybackslash}p{5.7cm}c>{\raggedright\arraybackslash}X}
\lsptoprule
{\bfseries Athematic Nasal-Infix Tr.} & & {\bfseries Thematic Nasal-Infix Intr.} \\
\midrule
{ \textit{*munék-\textsuperscript{ti}}/\textit{munk-\textsuperscript{tói}} \iw{PIE}{\textit{*munék}-/\textit{munk}-} ‘release, set free’} &  → & { *\textit{munk-é-\textsuperscript{ti}}\textsuperscript{} \iw{PIE}{*\textit{munk-é/ó}-} ‘go free, escape’} \\
{ Ved.\ \iw{Sanskrit}{\textit{muñcáti/te}} 
\textit{muñcáti}, Lat.\ \iwi{Latin}{\textit{ēmungō}}, -\textit{ere}} & & 
{ OCS \textit{promŭknǫ, -mŭknǫti sę, -mŭkŭ}\iw{OCS}{\textit{promŭknǫti sę, -mŭknǫ, -mŭkŭ}} ‘go out, spread', Lith.\ \textit{muñka, mùkti, mùko}\iw{Lith}{\textit{mùkti, muñka, mùko}} ‘break free’}\\
{ \textit{*linép-\textsuperscript{ti}}\textit{/linp-\textsuperscript{tói}} \iw{PIE}{\textit{*linép}-/\textit{linp}-} ‘make adhere to’} & →& { *\textit{linp-é-\textsuperscript{ti}} \iw{PIE}{*\textit{linp-é/ó}-} ‘cling, stick to’}\\
{ Ved.\ \iwi{Sanskrit}{\textit{limpáti}} AV+,\textsuperscript{?}Gk.\ \iwi{AG}{λιπαίνω}  ‘anoint’ Att.\ Ion.}  & & 
{ OCS \textit{prilĭ(p)nǫ, -lĭ(p)nǫti, -lĭpu}\iw{OCS}{\textit{prilĭ(p)nǫti, -lĭ(p)nǫ, -lĭpŭ}} ‘id.', Lith.\ \textit{lim̃pa, lìpti, lìpo}\iw{Lith}{\textit{lìpti}, \textit{lim̃pa}, \textit{lìpo}} ‘cling, stick to, climb', Go.\ \textit{aflifnan}\iw{Goth}{\textit{aflifnan, -lifniþ, -lifnoda}} ‘remain’}\\
\lspbottomrule
\end{tabularx}
\end{table}


In terms of its deeper PIE origins, \citet[200--214]{Gorbachov2007} argues that the formation is to be traced to a *\textit{h\textsubscript{2}}\textit{e}-conjugated version of the athematic nasal-infix present,\is{nasal present} which ultimately served to make \isi{anticausative}s and which as all *\textit{h\textsubscript{2}}\textit{e}-conjugation active present types underwent analogical \isi{thematization} on the way to NPIE.\footnote{See \citealt{Jasanoff2003} for the PIE *\textit{h\textsubscript{2}}\textit{e}-conjugation and its systematic \isi{thematization} in the Nuclear IE languages.} In my view, Gorbachov is correct in his analysis and reconstruction of this type for late PIE. The intention of this paper is to further support his analysis by outlining additional reflexes of the type in Ancient Greek and Indo-Iranian.\footnote{A further important feature of these presents is that in Slavic -- the only \isi{Northern Indo-European} branch to preserve the inherited IE aorist -- they systematically correlate with the thematic\is{thematic!aorist} aorist.}

\section{Ancient Greek}

In Ancient Greek (and further Classical Armenian\footnote{\label{fn:ClassArm}The Greek type has a close match in Classical Armenian, where, e.g., presents in \iwi{Armenian}{-\textit{ane/im}} regularly correlate with nasal-infix presents \is{nasal present} elsewhere, are made to thematic preterites, both inherited thematic aorists\is{thematic!aorist} and recategorized thematic imperfects, and show up in both \isi{transitive} and \isi{intransitive} formations. For the facts and discussion, see \citealt[162ff.]{Kocharov2019} with lit. The details of this type, whose core likely depends on a prehistoric Greek and Armenian genetic or convergence-related innovation, I leave for discussion elsewhere. For a useful orientational assessment of the relationship between Greek and Armenian, see most recently \citealt{Kim2018} with lit.}), it is possible to identify several sets of thematic nasal-affix presents \is{nasal present} that can be aligned with the \isi{Northern Indo-European} type outlined by Gorbachov -- cf.\ the five types listed in examples (\ref{ex:nu})--(\ref{ex:anw}).\footnote{Superscript \textsuperscript{A} marks verbs that at least optionally take an acc.\ object. Verbs that show strong Caland system\is{Caland!system (CS)}  associations, which are frequent across these categories, are marked with superscript \textsuperscript{C}. Note that all Caland-system-associative verbs\is{Caland!system (CS)} can in principle be understood as deadjectival \is{deadjectival!verb} provided they continue to stand in a synchronically transparent relationship with property concept\is{property concept (PC)} adjectives derived from the same roots. For primary verbs\is{derivation!deradical} as a foundational element of the PIE and IE Caland system,\is{Caland!system (CS)}  see \citealt{Rau2009}, \citeyear{Rau2013}, \citeyear{Rau2016}, \citeyear{Rau2017}, and further \citealt{Nussbaum2022}.} 


\begin{exe}
\ex \label{ex:nu} \textit{Nu}-presents,\footnote{The poetic verb \iwi{AG}{ἄνω}, -oμαι ‘accomplish, finish’ Hom.+ (: \iwi{AG}{ἁνύω}  ‘id.’ Hom.+) may also belong in this class if it is a relatively late innovation of the sort noted below. Another possibility is that the form is a recategorized thematic aorist,\is{thematic!aorist} either an inherited formation -- viz.\ \iwi{Sanskrit}{\textit{ásanat}} RV+ (: \iwi{Sanskrit}{\textit{asāniṣam}} RV): \iwi{Sanskrit}{\textit{sanóti}} RV, OAv.\  subj.\  \iwi{Avestan}{\textit{hanāt̰}} -- or an inner-Greek remodeling of a root aorist,\is{root!aorist} assuming the root is \iwi{PIE}{*\textit{senh\textsubscript{3}}-}. This latter analysis may have support in the peculiar aorist \iwi{AG}{ἤνεσα} \textit{IG} 7.3226, Orchomenos and elsewhere, if it represents the contamination product of this thematic aorist\is{thematic!aorist} and the \textit{s}-aorist otherwise well attested for this verb.}  e.g., 
\iwi{AG}{φθῑ́νω/φθῐ́νω}  ‘decay, waste away, disappear’ Hom.+: φθινύθω  Hom.+ (tr., intr.), ἔφθιτo Hom.+, ἔφθιoν Hom.+, ἔφθ(ε)ισα Hom.+, ἔφθιται Hom.+ 
(: Ved.\ \iwi{Sanskrit}{\textit{kṣiṇ{ā́}ti}}, \iwi{Sanskrit}{\textit{kṣiṇóti}} ‘destroy’ AV, YAv.\ \textit{jinā-\textsuperscript{ti}} \iw{Avestan}{\textit{jinā}-} ‘id.’; ON \textit{dvena} ‘disappear, decrease', ON \textit{dvína} \iw{ON}{\textit{dvena}, \textit{dvína}} ‘id.’)\footnote{As comparanda, I provide forms that are directly cognate and those that are closely parallel.}
\end{exe}




\begin{exe} \ex Nasal-infix to sonorant-plus-laryngeal-final roots
\begin{xlist}
    \ex \iwi{AG}{θάλλω}  ‘blossom, flourish’ Hes.\ \textit{hHom}.+: ἔθαλoν \textit{hHom}.+, τέθηλα \linebreak[5] Hom.+, \iwi{AG}{θαλέθω}  Hom.+ (: \iwi{AG}{θαλερός}  Hom.+, ἡ \iwi{AG}{θάλεια} Hom.+, τὸ \iwi{AG}{θάλoς}  Hom.+)\textsuperscript{C} (: Alb.\ 1sg.\ \iwi{Albanian}{\textit{dal}, \textit{del}} ‘go out’ < *\textit{dalnō}, \textit{dalnet} and MW \iwi{MWelsh}{\textit{deillyaw}*} {<}{<} *\textit{daln}-, see \citealt[257--259]{Schumacher2004} and  \citealt[968--969]{SchumacherMatzinger2013})
    \ex 
{\iwi{AG}{δῡ́νω}  ‘go into, enter; go down, sink’ Hom.+\textsuperscript{A}: \iwi{AG}{δύω, -oμαι}  Hom.+, ἔδῡν Hom.+ (intr.), ἔδῡσα, -αμην Hom.+}

\ex  \iwi{AG}{θῡ́νω}  ‘rush, dart along’ (: analogical \iwi{AG}{θῡνέω}  Hes.\ \textit{Sc}.; \iwi{AG}{θῡ́ω/θυίω}  ‘rage, seethe’ Hom.+, ἔθῡσα Call.) (: Ved.\ \iwi{Sanskrit}{\textit{dhūnóti}} ‘shake’ to \iwi{PIE}{*\textit{dʰeu̯H}-} [plus \iwi{AG}{θῡ́ω/θυίω}  ‘rage, seethe’ Hom.+, so \citetalias{Rix2001}], but \citetalias{Rix2001} \textit{s.r}. to \iwi{PIE}{*\textit{dʰeu̯}-} because of \iwi{AG}{θῡνέω}  Hes.\ \textit{Sc.}[!])\footnote{Given the structure of the root, the form is best assigned, \textit{pace} \citetalias{Rix2001}, who align it with the \textit{nu}-class above, to the category of nasal-infix formations to laryngeal final roots, with Ved.\ \iwi{Sanskrit}{\textit{dhūnóti}} an analogical innovation.}

\ex \iwi{AG}{κάμνω}  ‘(tr.) produce/acquire through labor/toil (Hom.[+], but rare); (intr.) labor, toil, grow weary toiling’ Hom.+\textsuperscript{A}: ἔκαμoν Hom+, κέκμηκα Hom.+ (: \iwi{AG}{ἀκάμας} , -αντoς  Hom.+; \iwi{Sanskrit}{\textit{śamnīṣe}} ‘labor’ YV, \iwi{Sanskrit}{\textit{śamāyáte}} ‘id.’ RV, \iwi{Sanskrit}{\textit{aśamīt}} ‘ist ruhig geworden’ AVP, \iwi{Sanskrit}{\textit{áśamanta}} ‘labored’ YV)

\ex ὀφείλω  Att.\ Ion., ὀφέλλω  Hom.\ Lesb.\ Arc., ὀφήλω  Cret.\ Arc.\ (Myc.) \iw{AG}{ὀφείλω, ὀφέλλω, ὀφήλω} ‘owe, be obliged to pay for/render; be bound, obliged to (+ inf.)’\textsuperscript{A}: ὦφλoν Hom.+ (Myc.), ὤφληκα Hom.+ Dial. (: \iwi{AG}{ὀφλισκάνω}  Att., etc.; τὸ \iwi{AG}{ὄφελoς}  Hom.+ [Myc.], \iwi{AG}{ὀφλητής}  [gloss]; cf.\ \citetalias{LIVaddenda}, *\textit{h\textsubscript{3}bʰelh\textsubscript{1}-})
\end{xlist}
\end{exe}


\begin{exe}
    \ex Nasal-infix to obstruent-plus-laryngeal-final roots
    \begin{xlist}
    \ex    aor.\ ($\Lleftarrow$ impf.) \iwi{AG}{ἔχανoν} ‘yawn, gape’ Hom.+: \iwi{AG}{κέχηνα} Hom.+, \iwi{AG}{χάσκω}  Att.\ Ion., \iwi{AG}{χαίνω}  late\textsuperscript{C}  (: OCS \iwi{OCS}{\textit{zinǫti}} ‘gape', OHG \iwi{OHG}{\textit{ginēn}} ‘yawn', ON \iwi{ON}{\textit{gína}} ‘id.’; \citetalias{LIVaddenda}, \iwi{PIE}{*\textit{g̑ʰeh\textsubscript{2}}-} and \textit{i}-present \iwi{PIE}{*\textit{g̑ʰeh\textsubscript{2}i}-})\footnote{Here I assume, with \citetalias{Rix2001}, that aor.\ \iwi{AG}{ἔχανoν} is an aspectually recategorized imperfect. A parallel shift is responsible for \isi{perfective} aspect of OCS \iwi{OCS}{\textit{zinǫti}}.}

\ex \iwi{AG}{φθάνω/φθᾱ́νω}  ‘(tr.) be beforehand,  anticipate; (intr.) come/act first ($\pm$part.)’ Hom.+\textsuperscript{A}: ἔφθην Hom.+ Att., φθάμενoς  Hom.\ Hes., ἔφθασα Att.\ Ion.\  (: Ved.\ opt. \iwi{Sanskrit}{\textit{daghnuyāt}} ‘würde knapp verfehlen’ KS, \textit{m{ā́} dhak} ‘soll nicht verfehlen’ RV. See \citetalias{Rix2001} \iwi{PIE}{*\textit{dʰeg\textsuperscript{u̯h}}\textit{h\textsubscript{2}}{}-} for further possible comparanda.)
    \end{xlist}

\end{exe} 


\begin{exe} 
\ex Nasal-infix to stop-final root
\begin{xlist}
    \ex \iwi{AG}{λάμπω}  ‘shine’ Hom.+ (: \iwi{AG}{λαμπρός}  Hom.+;  Hitt. \textit{lap-\textsuperscript{zi}}, \iw{Hitt}{\textit{lap}-} pret.\ \textit{lāp-\textsuperscript{ta}} [cf. Lith.\ \iwi{Lith}{\textit{lóp\.e}}, etc.])\textsuperscript{C}
\ex  {[ὁ \iwi{AG}{θρόμβoς}  ‘lump, clot’ Att.\ Ion.: \iwi{AG}{τρέφω, -ομαι} ‘make firm, congeal’, ἔτραφoν Hom.\ (intr.), ἔτρεψα, τέτρoφα Hom.+\textsuperscript{C} 
 (: \iwi{AG}{ταρφύς}  Hom.+,\linebreak[4] τάρφα Hom.+, τὸ \iwi{AG}{τάρφoς}  Hom.+; Lith.\ \textit{drim̃ba, drìbti}\iw{Lith}{\textit{drìbti, drim̃ba}} ‘fall in flakes, drop’)]}\footnote{This and the following two nouns obviously depend on an underlying nasal-infix present,\is{nasal present} either an athematic \isi{transitive} or a thematic \isi{intransitive} of the sort under discussion. I assign them to the latter type, first because of the semantic profile of the roots and their derivational \textit{averbo}, which -- at least in the case of ὁ \iwi{AG}{θρόμβoς}  and τὸ \iwi{AG}{θάμβoς}  -- is exactly parallel to the other verbs that can be assigned to this class, and second, because the retention of the nasal-infix formant  in a nominal derivative is more easily explained as the result of a synchronically isolated thematic nasal-infix \isi{intransitive} \is{nasal present} -- replaced prehistorically in Greek by the -ανω  type below -- than an otherwise reasonably well-paralleled athematic nasal-infix \isi{transitive}.} %Cf. YAv.\ nom. sg.\ \textit{θra̜fš}ᵒ ‘satisfaction (\textit{vel sim}.)’ Yt.5.26, YAv.\ \textit{θra̜fəδra}- ‘rich in; satisfied’ and synchronically isolated Ved.\ \textit{tr̥mpáti} ‘take delight in, enjoy (+gen.)’. RV cited in Section \ref{sec:IIr} below.}
\ex {[τὸ \iwi{AG}{θάμβoς}  ‘astonishment’ Hom.+: \iwi{AG}{θαμβέω}  ‘be astonished’ (: ἐθάμβησε Hom.+)\textsuperscript{A}, ἔταφoν Hom.+, τέθηπα Hom.+ (: \iwi{AG}{θώψ}, θωπός  Att.\ Ion.\textsuperscript{C};  deverbal\textsuperscript{?}\is{deverbal/deverbative!verb} PGmc.\ \iwi{PGmc}{*\textit{dumba}-}). Non-nasal root, \textit{pace} \citealt[237]{Hackstein2002Epen} with lit.\footnote{The long vowel in the \isi{perfect} \iwi{AG}{τέθηπα} Hom.+ and the root noun\is{root!noun} \iwi{AG}{θώψ}, θωπός  Att.\ Ion., whose secondary origins\is{derivation!stem-based} are difficult to motivate, make the reconstruction of an underlying non-nasal root the default assumption.}]}
\ex {[ὁ \iwi{AG}{στρόμβoς}  ‘top; spinning object’ Hom.+: \iwi{AG}{στρέφω, -oμαι} ‘twist, turn’ Hom.+, ἔστρεψα, -αμην Hom.+]} Cf.\ YAv.\ nom.\ sg.\ \iwi{Avestan}{\textit{θra̜fš}°} ‘satisfaction (\textit{vel sim}.)’ Yt.5.26, YAv.\ \iwi{Avestan}{\textit{θra̜fəδra}-} ‘rich in; satisfied’: Ved.\ \iwi{Sanskrit}{\textit{tr̥mpáti}} ‘take delight in, enjoy (+gen.)’ RV. See Section \ref{sec:IIr}.
\end{xlist}
\end{exe}


\begin{exe} 
\ex \label{ex:anw} \iwi{AG}{-άνω}  to obstruent final roots\footnote{This is a productive\is{productivity} present-stem formant in Greek, which is regularly used to generate presents from thematic aorist\is{thematic!aorist} forms. I mostly omit obviously -- e.g., \iwi{AG}{ἀλφάνω}  Att.\ (: \iwi{AG}{ἄλφω}  Hsch.): ἦλφoν Hom.+ -- and likely -- e.g., \iwi{AG}{λιμπάνω}  Hom.\ (\textit{v.l}.) Sa. Th. Hp. (+) (: \iwi{AG}{λείπω}  Hom.+): ἔλιπoν Hom.+ -- secondary creations. I include some obviously or likely secondary, inner-Greek formations here to illustrate more fully the semantic subcategories and coherence of this subtype. I use \textsuperscript{\$} to mark such formations as secondary, Greek-internal innovations. Forms are arranged according to semantic subtype, from property concept\is{property concept (PC)} to perception/cognition and possession, \textit{inter alia}.}
\begin{xlist}
    \ex { \iwi{AG}{ἁνδάνω}  ‘be pleasing’ Hom.+ (: \iwi{AG}{ἥδoμαι} Att.\ Ion.\ Dial.): ἕαδoν Hom.+, ἕαδα Hom.+ (: \iwi{AG}{ἡδύς}  Hom.+, τὸ \iwi{AG}{ἧδoς}  Hom.+, etc.\textsuperscript{C}; Ved.\ \iw{Sanskrit}{\textit{s\textsubscript{(u)}vádati/te}} \textit{s\textsubscript{(u)}vádati} ‘make sweet, sweeten’ RV, \textit{sv{ā́}date}\iw{Sanskrit}{\textit{svā̆́date}} ‘become sweet’ RV)} 

\ex { \iwi{AG}{ἐκφλυνδάνει} ‘break out (of sores)’ Hp. (: \iwi{AG}{φλυδαρός}  ‘soft, flabby’ Hp. \textit{ap}. Gal.19.152, \iwi{AG}{φλυδάω}  ‘have excess of moisture, be soft/flabby’ Hp.)\textsuperscript{C}} 

\ex { \iwi{AG}{ἀργάνω}  ‘be white’ E.+ (: \iwi{AG}{ἀργός}  Hom.+, τὸ \iwi{AG}{ἄργoς}  Hom.+, etc.)\textsuperscript{\$C}}

\ex \iwi{AG}{oἰδάνω}  ‘swell (intr.)’ I646 (oἰδάνεται), Ar.\textit{Pax}.116 (oἰδάνων); ‘make swell’ I554 ({$\simeq$} A.R.1.478): \iwi{AG}{oἰδαίνω}  ‘swell (intr.)’ A.R.+ (tr., late and rare), \iwi{AG}{oἰδέω}  ‘id.’ Hom.+ (: \iwi{AG}{oἰδαλέoς}  Hom.+, τὸ \iwi{AG}{oἶδoς}  Hom.+, etc.\textsuperscript{C}; ClArm. \iwi{Armenian}{\textit{aytnowm}} ‘swell [intr. \textit{tantum}]’)\footnote{\label{fn:itract}Intransitive\is{intransitive} oἰδάνων at Ar.\textit{Pax}.116 is an obvious \textit{forma difficilior} and points to an original \isi{intransitive} meaning for the active. The Homeric opposition \isi{transitive} active: \isi{intransitive} middle has been analogically created according to a well-attested Greek morphological restructuring process -- cf., e.g., {αἴθω} \iw{AG}{αἴθω, -oμαι} ‘burn (intr.)’ Pi.S., αἴθoυσα ‘portico’ Hom.+: αἴθω  ‘burn (tr.)', {αἴθoμαι} \iw{AG}{αἴθω, -oμαι} ‘burn (intr.)’ Hom.+. The \isi{intransitive} active is also supported by largely \isi{intransitive} and synchronically/semantically isolated \iwi{AG}{oἰδαίνω}  ‘swell (intr.)’ A.R.+ (tr., late and rare), which is likely a \isi{backformation} from an innovative aorist to present \iwi{AG}{oἰδάνω}. Similar considerations hold for immediately following \isi{transitive} and \isi{intransitive} active \iwi{AG}{κῡδάνω}  ‘be triumphant (intr., Y.42); glorify (tr., ξ73).’} 

\ex { \iwi{AG}{κῡδάνω}  ‘be triumphant (intr., Y.42); glorify (tr., ξ73)’: \iwi{AG}{κῡδαίνω}  ‘glorify, honor (tr. \textit{tantum})’ Hom.+, ἐκῡ́δηνα Hom.+ (: \iwi{AG}{κῡδρός}  Hom.+, \iwi{AG}{κῡδάλιμoς}  Hom.+, τὸ \iwi{AG}{κῦδoς}  Hom.+, etc.)\textsuperscript{C}}

\ex { \iwi{AG}{λανθάνω, -oμαι} ‘escape notice ($\pm$acc.); forget’ (: \iwi{AG}{λήθω, -oμαι}  Hom.+)\textsuperscript{A}: ἔλαθoν, -oμην Hom.+: λέληθα Hom.+ (: \iwi{AG}{λάθρηι} Hom.+, etc.)\textsuperscript{\$C}}

\ex { \iwi{AG}{ὀλισθάνω}  ‘slip’ Att.\ (: \iwi{AG}{ὀλισθαίνω}  ‘id.’ Arist.+): ὤλισθε Hom.+ (: Go.\ \iwi{Goth}{\textit{slindan}} ‘schlingen’; Ved.\ \iwi{Sanskrit}{\textit{srédhati}} ‘gleitet aus, handelt falsch’ RV, \iwi{Sanskrit}{\textit{sridhat}} RV)\textsuperscript{C}} 

\ex { \iwi{AG}{ἁμαρτάνω}  ‘miss the mark, go astray’ Hom.+: ἥμαρτoν Hom.+} 

\ex { \iwi{AG}{τυγχάνω}  ‘chance on, happen, hit the mark’ Hom.+: ἔτυχoν Hom.+, τετύχηκα Hom.+ (: \iwi{AG}{τεύχω}  ‘make ready, produce’ Hom.+, ἔτευξα Hom.+, τετευχώς  Hom.\ [Myc.])}

\ex { \iwi{AG}{πυνθάνoμαι} ‘learn, ascertain’ Hom.+\textsuperscript{A} (: \iwi{AG}{πεύθoμαι} Hom.+): ἐπυθόμην Hom.+: πέπυσμαι Hom.+ (: OCS \textit{vŭzbŭnǫti, vŭzbŭde}\iw{OCS}{\textit{vŭzbŭ(d)nǫti, -bŭ(d)nǫ, -bŭdŭ}} ‘wake up', Lith.\ \textit{atbuñda, atbùsti, atbùdo}\iw{Lith}{\textit{atbùsti, atbuñda, atbùdo}} ‘id.’; OIr.\ \iwi{OIr}{\textit{ad}$\cdot$\textit{boind}})\textsuperscript{C}}

\ex \iwi{AG}{αἰσθάνoμαι} ‘perceive, take notice of’ Att.\ Ion.\textsuperscript{(A)}: ηἰσθόμην Att.\ Ion., \iwi{AG}{ηἴσθημαι} Att.\ Ion.\textsuperscript{\$}

\ex  \iwi{AG}{μανθάνω}  ‘learn, notice, understand’ Pi+\textsuperscript{A}: ἔμαθoν Hom.+, μεμάθηκα Att.\ Ion.\textsuperscript{\$}

\ex  \iwi{AG}{δαρθάνω}  ‘fall asleep’ Hierocl.: ἔδαρθoν Hom.+\textsuperscript{\$}

\ex  \iwi{AG}{θιγγάνω}  ‘touch, handle’ Att.\ Ion.+: ἔθιγoν Archil.+\footnote{On this verb, see the recent discussion of \citet{Nikolaev2023}.}

\ex  \iwi{AG}{χανδάνω}  ‘take hold; be capable’ Hom.+\textsuperscript{A}: ἔχαδoν Hom.+, κέχα/oνδα Hom.\  (: Lat.\ \iwi{Latin}{\textit{prehendō}}, -\textit{ere} ‘seize, grasp’, OIr.\ \iwi{OIr}{\textit{ro}$\cdot$\textit{geinn}} ‘occupies’, Go.\ \iwi{Goth}{\textit{duginnan}} ‘begin’\footnote{See \citet{Fortson2009}, and further \citet{Jasanoff2022}, who proposes a special origin for this set of verbs and at least this subset of -ανω  presents \is{nasal present} in Greek (and Armenian). In my view, this set is best understood as an instance of the general type under discussion, with Lat.\ \iwi{Latin}{\textit{prehendō}} as the continuation of \iwi{PIE}{*\textit{gʰn̥dʰ-e/o}-}, Go.\ \iwi{Goth}{\textit{duginnan}} as a re-characterization of this inherited stem according to the productive\is{productivity} Germanic type, and OIr.\ \iwi{OIr}{\textit{ro}$\cdot$\textit{geinn}} ‘occupies’ as the phonologically regular development of the same stem \textit{à la} \citet{Fortson2009}.})

\ex  \iwi{AG}{λαγχάνω}  ‘obtain by lot, come to possess’ Hom.+\textsuperscript{(A)}, intr. ‘fall to one’s lot’ Hom.+: ἔλαχoν Hom.+, λέλoγχα Hom.

\ex  \iwi{AG}{λαμβάνω}  ‘take hold, grasp, seize’ Att.\ Ion.\textsuperscript{A} (: \iwi{AG}{λάζoμαι} Hom., Dial.): ἔλαβoν Hom.+, εἴληφα Att./λελάβηκα Ion., Dial.\textsuperscript{\$} 
\end{xlist}
\end{exe}




It is possible to divide this material into two basic sets of forms. The first set, which includes everything but the final άνω-class in ex.\ (\ref{ex:anw}), consists of verbs that are phonologically directly derivable from a thematic nasal-infix\is{nasal present} formation of the Gorbachov type made to the PIE root in question, with a subtype (a) that pairs with a root aorist\is{root!aorist} or for which the aorist is not attested -- viz.\ \iwi{AG}{φθῑ́νω/φθῐ́νω}  Hom.+: ἔφθιτo Hom.+, ἔφθιoν Hom.+, ἔφθ(ε)ισα Hom.+, \iwi{AG}{δῡ́νω}  Hom.+\textsuperscript{A}: ἔδῡν Hom.+ (intr.), ἔδυσα, -αμην Hom.+, etc. -- and a subtype (b) that pairs with a thematic aorist\is{thematic!aorist} -- viz.\ \iwi{AG}{θάλλω}  Hes.\ \textit{hHom}.+: ἔθαλoν \textit{hHom}.+, etc. The second set consists of verbs in \iwi{AG}{-άνω}, a highly productive\is{productivity} and expansive category in both prehistoric and historical Greek. As far as can be demonstrated, the verbs found in this second set are not made to roots in a final laryngeal and so cannot at first sight be explained as the direct phonological results of the thematic nasal-infix formation in question -- many if not all must be of analogical origin. Nevertheless, it is important to note that these verbs largely show the same semantics as set 1 (see further below ),  systematically pair with the thematic aorist\is{thematic!aorist} (viz.\ \iwi{AG}{ἁνδάνω}  Hom.+: ἕαδoν Hom.+, etc.), as set 1 (and further as in Slavic), have a substantial number of forms that correlate with the Caland system,\is{Caland!system (CS)}  as set 1, and of course share a similar thematic active nasal stem formant, exactly as set 1. There is thus every reason to assume that both sets continue the same underlying type.

Across both sets of forms, these verbs regularly correlate with \isi{nasal present}s in other IE languages. Some have exact or close formal and semantic matches of the type that Gorbachov outlined, while others have athematic nasal-infix cognates \is{nasal present} where they pattern semantically with the \isi{intransitive} middle of the athematic nasal-infix cognate but not the \isi{transitive} active, cf.\ (\ref{ex:semparallel})--(\ref{ex:mismatch}).

\begin{exe}
    \ex \label{ex:semparallel} Close semantic and morphological equation or parallel\footnote{A likely further instance of this type underlies ὁ θρόμβoς  ‘lump, clot’ Att.\ Ion.: Lith.\ \textit{drim̃ba, drìbti}\iw{Lith}{\textit{drìbti, drim̃ba}} ‘fall in flakes, drop’.}
       \begin{xlist}
\ex { \iwi{AG}{φθῑ́νω/φθῐ́νω}  ‘disappear, decrease’ Hom.+: ON \textit{dvena} ‘disappear, decrease', ON \textit{dvína} \iw{ON}{\textit{dvena}, \textit{dvína}} ‘id.’ (: Ved.\ \iwi{Sanskrit}{\textit{kṣiṇ{ā́}ti}}, \textit{kṣiṇóti} AV, YAv.\ \textit{jinā-\textsuperscript{ti}}) \iw{Avestan}{\textit{jinā}-} }
\ex { \iwi{AG}{θάλλω}  ‘blossom, flourish’ Hes.\ \textit{hHom}.+: MW \iwi{MWelsh}{\textit{deillyaw}*}, Alb.\ \iwi{Albanian}{\textit{dal}, \textit{del}} ‘go out’}

\ex { \iwi{AG}{ἔχανoν} ‘gape, yawn’ Hom.+: OCS \iwi{OCS}{\textit{zinǫti}} ‘gape', OHG \iwi{OHG}{\textit{ginēn}} ‘yawn’, ON \iwi{ON}{\textit{gína}} ‘id.’}

\ex { \iwi{AG}{πυνθάνoμαι} ‘learn, ascertain’ Hom.+: OCS \textit{vŭzbŭnǫti, vŭzbŭde}\iw{OCS}{\textit{vŭzbŭ(d)nǫti, -bŭ(d)nǫ, -bŭdŭ}} ‘wake up', Lith.\ \textit{atbuñda, atbùsti, atbùdo}\iw{Lith}{\textit{atbùsti, atbuñda, atbùdo}} ‘id.’ (: OIr.\ \iwi{OIr}{\textit{ad}$\cdot$\textit{boind}})}

\ex { \textsuperscript{?}\iwi{AG}{χανδάνω}  ‘take hold’ Hom.+: Lat.\ \iwi{Latin}{\textit{prehendō}}, -\textit{ere}, OIr.\ \iwi{OIr}{\textit{ro}$\cdot$\textit{geinn}} `occupies’, Go.\  
 \iwi{Goth}{\textit{duginnan}} ‘begin’ (cf.\ \citealt{Fortson2009}; \citealt{Jasanoff2022})}
    \end{xlist}
\ex Semantic correspondence with athematic nasal-infix middle
    \begin{xlist}
        \ex   { \iwi{AG}{κάμνω}  ‘(tr.) produce/acquire through labor/toil (Hom.[+], but rare); (intr.) labor, toil, grow weary laboring/toiling’ Hom.+: \iwi{Sanskrit}{\textit{śamnīṣe}} ‘labor’ YV, \iwi{Sanskrit}{\textit{śamāyáte}} ‘id.’ RV, \iwi{Sanskrit}{\textit{aśamīt}} ‘ist ruhig geworden’ AVP, \iwi{Sanskrit}{\textit{áśamanta}} ‘labored’ YV}
\ex { \iwi{AG}{ἁνδάνω}  ‘be/become pleasing’ Hom.+: \textsuperscript{?}Ved.\ \iw{Sanskrit}{\textit{s\textsubscript{(u)}vádati/te}} \textit{s\textsubscript{(u)}vádate} ‘wird schmackhaft’ RV\footnote{Assuming this form continues a nasal-infix present.\is{nasal present}} (: Ved.\ \iw{Sanskrit}{\textit{s\textsubscript{(u)}vádati/te}} \textit{s\textsubscript{(u)}vádati} ‘macht schmackhaft’ RV)}
\ex { \iwi{AG}{oἰδάνω}  ‘swell (intr.)’ I646 (oἰδάνεται), Ar.\textit{Pax}.116 (oἰδάνων); ‘make swell’ (tr.) I554 ({$\simeq$} A.R.1.478): ClArm. \iwi{Armenian}{\textit{aytnowm}} ‘swell [intr.\ \textit{tantum}]’} 
    \end{xlist}
    \ex \label{ex:mismatch} Semantic mismatch with athematic nasal-infix \is{nasal present} active\footnote{For the pair \iwi{AG}{φθάνω/φθᾱ́νω}  ‘(tr.) be beforehand, outstrip, anticipate; (intr.) come/act first ($\pm$part.)’ Hom.+: Ved.\ opt. \iwi{Sanskrit}{\textit{daghnuyāt}} ‘würde knapp verfehlen', it is important to note that active \textit{nu}-presents \is{nasal present} are often \isi{intransitive} in Vedic and in some cases replace the thematic nasal-infix \isi{intransitive} actives discussed here -- viz.\ Ved.\ \iwi{Sanskrit}{\textit{tr̥mpáti}} ‘take delight in’ RV: \iwi{Sanskrit}{\textit{tr̥pṇóti}} RV ‘id.’}
    \begin{xlist}
        \ex { \iwi{AG}{φθῑ́νω/φθῐ́νω}  ‘disappear, decrease’ Hom.+: Ved.\ \iwi{Sanskrit}{\textit{kṣiṇ{ā́}ti}}, \iwi{Sanskrit}{\textit{kṣiṇóti}} ‘destroy’ AV, YAv.\ \textit{jinā-\textsuperscript{ti}} \iw{Avestan}{\textit{jinā}-} ‘id.’ (: ON \textit{dvena} ‘disappear', decrease', ON \textit{dvína} \iw{ON}{\textit{dvena}, \textit{dvína}} ‘id.’)}
\ex { \iwi{AG}{θῡ́νω}  ‘rush, dart along’: Ved.\ \iwi{Sanskrit}{\textit{dhūnóti}} ‘shake’} 
\ex { \iwi{AG}{ἁνδάνω}  ‘be/become pleasing’ Hom.+: \textsuperscript{?}Ved.\ \iw{Sanskrit}{\textit{s\textsubscript{(u)}vádati/te}}  \textit{s\textsubscript{(u)}vádati} ‘macht schmackhaft’ RV, \iwi{Sanskrit}{\textit{svā̆́date}} ‘wird schmackhaft’ RV}
\ex { \iwi{AG}{ὀλισθάνω}  ‘slip’ Att.\ (: \iwi{AG}{ὀλισθαίνω}  ‘id.’ Arist.+): ὤλισθε Hom.+: Go.\ \iwi{Goth}{\textit{slindan}} ‘schlingen’ (: Ved.\ \iwi{Sanskrit}{\textit{srédhati}} RV, \textit{sridhat} RV)} 
\ex { \iwi{AG}{πυνθάνoμαι} ‘learn, ascertain’ Hom.+ (: \iwi{AG}{πεύθoμαι} Hom.+):  OIr.\ \iwi{OIr}{\textit{ad}$\cdot$\textit{boind}} (: OCS \textit{vŭzbŭnǫti, vŭzbŭde}\iw{OCS}{\textit{vŭzbŭ(d)nǫti, -bŭ(d)nǫ, -bŭdŭ}} ‘wake up', Lith.\ \textit{atbuñda, atbùsti, atbùdo}\iw{Lith}{\textit{atbùsti, atbuñda, atbùdo}} ‘id.’)}
    \end{xlist}
 
\end{exe}


 


The semantic patterning here aligns closely with the fact that nearly two-thirds of the forms in Greek are \isi{intransitive} and that \isi{transitive} verbs are largely limited to cases where an original accusative of respect or goal is likely, such as λανθάνω, \mbox{-oμαι}\iw{AG}{λανθάνω, -oμαι} ‘escape notice ($\pm$acc.); (mid.) forget’ (: \iwi{AG}{λήθω, -oμαι} Hom.+) or \iwi{AG}{φθάνω/φθᾱ́νω} ‘(tr.) be beforehand, outstrip, anticipate; (intr.) come/act first ($\pm$part.)’ Hom.+, or where the root denotes a cognitive-perceptual or possessional state, state-oriented semantic categories that are regularly \isi{transitive} cross-linguistically -- e.g., \iwi{AG}{πυνθάνoμαι} ‘learn, ascertain’ Hom.+ (: \iwi{AG}{πεύθoμαι} Hom.+) or \iwi{AG}{χανδάνω}  ‘take hold of’ Hom.+.\footnote{On state-oriented semantic subclasses, see \citealt{Rau2013}, \citeyear{Rau2016}, and \citeyear{Rau2017}.}

The traditional approach to these verbs, which stretches back to the 19\textsuperscript{th} c. and is followed by all modern scholars, including  \citealt{Klingenschmitt1982}, \citealt{Rix1992},  \citetalias{Rix2001} \textit{passim}, and others, explains them as the result of an inner-Greek functionless \isi{thematization}. This is problematic on all fronts and can be excluded.\footnote{\label{fn:themat}The traditional approach, which explains the whole type this way and especially forms with suffixal \iwi{AG}{-άνω}  $\Lleftarrow$ -\textit{Cn̥h\textsubscript{x}ont(i)} \emph{bzw}.\ -\textit{Cn̥u̯ont(i)}, has four fatal issues:
\begin{enumerate}
    \item  As noted in the survey, the Greek forms show a fundamental and unmotivated semantic mismatch with active athematic nasal-infix presents,\is{nasal present} which are in effect exclusively \isi{transitive} in the languages and mostly make \isi{causative} \is{change of state (CoS)} change of state/location ($\pm$\isi{activity}) formations -- apart from late NPIE *\textit{dʰr̥snéu̯-\textsuperscript{ti}} \iw{PIE}{*\textit{dʰr̥snéu̯}-} (: Ved.\ \iwi{Sanskrit}{\textit{dhr̥ṣṇóti}} ‘dare', etc.) and some \textit{nu}-presents \is{nasal present} in Indic and Iranian, which are obviously secondary. On the general semantic profile of nasal-infix presents,\is{nasal present} see \citealt{Meiser1993}. 
    \item There is no reason to assume that Greek generalized the thematic 3\textsuperscript{rd} pl. in \iwi{PIE}{*-\textit{ont(i)}} to athematic stems and positive evidence against this assumption. See \citealt{Peters2004} and further \citealt[286--298]{Willi2018}.
    \item There is no reason to assume that the syllabification of the 3\textsuperscript{rd} pl. of the nasal-infix presents \is{nasal present} in Greek or the other IE languages was -\textit{Cn̥h\textsubscript{x}}\textit{ont(i)} \emph{bzw}.\ -\textit{Cn̥u̯ont(i)}. The nasal-infix formant fails to show syllabic variants in its \textit{inflectional} paradigm (see  \citealt{Praust2004}), at least in the environments in question here -- viz.\ when infixed before a sonorant or laryngeal in a minimally tri-consonantal root -- and instead developed an epenthetic vowel, a situation also evidenced by the Greek treatment of the nasal-infix presents \is{nasal present} to laryngeal final roots. The default assumption is that this also held in the 3\textsuperscript{rd} pl. and that the epenthetic vowel was later deleted as in Indo-Iranian.
  \item Even if these assumptions were reasonable, there is no other support for the assumption that nasal-infix presents \is{nasal present} were regularly or even occasionally subject to \isi{thematization} in the deeper prehistory of the language. Nasal-infix presents \is{nasal present} instead follow other well documented developmental pathways,  see \citealt{Sturm2021}. 
  \begin{enumerate}
      \item Some have replaced their athematic forms with \textit{gr̥bhāyáti}-derivatives -- e.g., \iwi{AG}{ὑφαίνω}  ‘weave’ (: Ved.\ \iwi{Sanskrit}{\textit{unábdhi}} ‘id.’), φαίνω\iw{AG}{φαίνω, -oμαι}  ‘show, make appear’ (: ClArm. \iwi{Armenian}{\textit{banam}} ‘disclose, discover’), \iwi{AG}{ἰαίνω}  ‘warm, cheer’ (possible athematic \iwi{AG}{ἰνάω}  ‘empty’) (: Ved.\ \iwi{Sanskrit}{\textit{iṣṇ{ā́}ti}} ‘set in vigorous motion’), \iwi{AG}{κλῑ́νω}  ‘incline’ (: YAv.\ -\textit{sirinao}-\textit{\textsuperscript{ti}} \iw{Avestan}{-\textit{sirinao}-} ‘lean on (tr.)’), \iwi{AG}{κρῑ́νω}  ‘judge’ (: Lat.\ \iwi{Latin}{\textit{cernō}}, -\textit{ere} ‘distinguish, discern’). On the type, see \citealt{Jasanoff2003}, and on the Greek forms, \citealt{Kolligan2010}.
      \item Nasal-infix stems to roots in a stop were regularly remodeled as \textit{nu}-stems -- cf.\ \iwi{AG}{ζεύγνῡμι} ‘yoke, join’ (: Ved.\ \iwi{Sanskrit}{\textit{yunákti}} ‘id.’), \iwi{AG}{ὀρέγνῡμι} ‘stretch out’ (: Ved.\ \iwi{Sanskrit}{\textit{r̥ṇákti}*} ‘id.’), etc.
      \item Nasal-infix stems to roots in \textit{{}-h}\textit{\textsubscript{2}} develop normally and then in Attic Ionic are restructured as \textit{nu}{}-verbs -- cf.\ \iwi{AG}{πίτνημι} ‘spread out’ Hom.+ (: \iwi{AG}{πετάννῡμι} Att.), etc. 
      \item Nasal-infix stems to roots in \textit{{}-h}\textit{\textsubscript{3}} develop normally and then in the history of Greek are restructured as \textit{nu}-stems, with further late remodeling in some cases -- cf.\ \iwi{AG}{στόρνῡμι} ‘strew’ Hom.+ (: \iwi{AG}{στoρέννῡμι} Eust.), etc. 
      \item Nasal-infix stems to roots in -\textit{h\textsubscript{1}} are early analogically restructured as thematic stems -- cf.\ \iwi{AG}{βάλλω}, \iwi{AG}{τάμνω}, etc.
  \end{enumerate}
\end{enumerate}  The nasal-infix stems that otherwise show up as thematic in Greek are 1) late \isi{backformation}s from the aorist -- cf., e.g., \iwi{AG}{τανύω}  ‘stretch’ Hom.+ (: τάνυμαι Hom.), \iwi{AG}{ἁνύω}  ‘seek, strive’ Hom.+ (: ἅνυμες  Theoc.+), 2) of denominative origin\is{denominative!verb} -- cf., e.g., \iwi{AG}{δῑνέω}  ‘whirl around', \iwi{AG}{κῑνέω}  ‘set in motion', on which see \citealt{Nikolaev2025MSS, Nikolaev2025glotta} -- or obviously secondary -- \iwi{AG}{ἱκνέoμαι} ‘reach, arrive’ Hom.+ Att.\ Ion.\ (: \iwi{AG}{ἵκω}  Hom.\ Dial.) --, 3) likely due to reinterpretation of archaic volitional/\isi{desiderative} subjunctives as indicatives in performative, speech-act contexts \textit{et sim}. -- viz.\ \iwi{AG}{ὀμνύω}  ‘swear’ Hom.+ Dial. (: \iwi{AG}{ὄμνῡμι} Hom.+), \iwi{AG}{τῑ́νω/τίνω}  ‘pay’ Hom.+ (: <τιννυνα(ι)> Kyme, ca.\ 700 BCE, 
 \citealt[201--214]{BartonekBuchner1995}, τίνυμαι Hom.\ Hes.\ Hdt. E.) --, or 4) simply morphologically obscure -- viz.\ \iwi{AG}{ἀρνέoμαι} ‘deny’ Hom.+, \iwi{AG}{κυνέω}  ‘kiss’ Hom.+.} 

The obvious alternative to this approach is to motivate the form and function of these verbs by aligning them with the active thematic nasal-infix \is{nasal present} \isi{inchoative} type that Gorbachov has outlined for Baltic, Slavic, and Germanic. This is straightforward in the case of the forms in set 1) above -- that is, everything but the άνω-presents.\is{nasal present}  These can all easily be explained as direct phonological and semantic continuants of this type -- viz.\ \iwi{AG}{φθῑ́νω/φθῐ́νω}  Hom.+ < \iwi{PIE}{\textit{*dʰg\textsuperscript{u̯h}inu̯-o/e}-} (: \iwi{AG}{φθινύθω}  Hom.+, Ved.\ \iwi{Sanskrit}{\textit{kṣiṇ{ā́}ti}}, \iwi{Sanskrit}{\textit{kṣiṇóti}} ‘destroy’ AV, YAv.\ \textit{jinā-\textsuperscript{ti}} \iw{Avestan}{\textit{jinā}-} ‘id.’; ON \textit{dvena} ‘disappear, decrease', ON \textit{dvína} \iw{ON}{\textit{dvena}, \textit{dvína}} ‘id.’), \iwi{AG}{θάλλω}  Hes.\ \textit{hHom}.+ < \iwi{PIE}{\textit{*dʰl̥nh\textsubscript{1}-o/e}-} (: Alb.\ 1sg.\ \iwi{Albanian}{\textit{dal}, \textit{del}} ‘go out’ < *\textit{dalnō}, \textit{dalnet} and MW \iwi{MWelsh}{\textit{deillyaw}*} {<}{<} *\textit{daln}-), etc. 

The forms in set 2) -- άνω-presents \is{nasal present} to roots in obstruent without a demonstrable laryngeal --, which largely show the same morphology, semantics, and derivational relations as set 1) and cannot be explained as the result of an inner-Greek functionless \isi{thematization} (see footnote \ref{fn:themat}), can also be aligned with this type under two straightforward assumptions. First, the suffix \iwi{AG}{-άνω}  was extracted from roots that ended in an obstruent plus laryngeal -- viz.\ -\textit{Cano/e}{}- < *-\textit{Cn̥h\textsubscript{x}}\textit{e}/\textit{o}{}- -- and then extended to all obstruent final roots. This development can be understood as an example of the “externalization” of the nasal formant, a process commonly found in all IE languages and well-attested in Greek -- cf., e.g., \iwi{AG}{ζεύγνῡμι} (: Ved.\ \iwi{Sanskrit}{\textit{yunákti}}), \iwi{AG}{ὀρέγνῡμι} (: Ved.\ \iwi{Sanskrit}{\textit{r̥ṇákti}*}), etc. or, with the \textit{gr̥bhāyáti}{}-type, Gk.\ \iwi{AG}{ὑφαίνω}  (: Ved.\ \iwi{Sanskrit}{\textit{unábdhi}}), \iwi{AG}{κῡδαίνω}  ‘glorify, honor’ Hom.+ ($\Lleftarrow$ \iwi{AG}{ἰαίνω}  ‘invigorate, warm’ [: Ved.\ \iwi{Sanskrit}{\textit{iṣṇ{ā́}ti}} ‘set in vigorous motion’]), etc. -- and can be seen to be exactly parallel to what has happened in Slavic and Germanic, where the synchronically normal shape of the formant has been extracted from sonorant-plus-laryngeal-final roots -- viz.\ -\textit{I/R̥ne/o{}-} < \textit{*-I/R̥nh\textsubscript{x}}\textit{e/o-} -- and extended at the expense of all other stem shapes -- e.g., OCS \textit{prilĭ(p)nǫ, -lĭ(p)nǫti}\iw{OCS}{\textit{prilĭ(p)nǫti, -lĭ(p)nǫ, -lĭpŭ}} ‘cleave, cling to', Go.\ \textit{aflifnan}\iw{Goth}{\textit{aflifnan, -lifniþ, -lifnoda}} ‘remain', ON \iwi{ON}{\textit{lifna}} ‘id.’ (: Lith.\ \textit{lim̃pa, lìpti}\iw{Lith}{\textit{lìpti}, \textit{lim̃pa}, \textit{lìpo}} ‘cling, stick to; climb’: Ved.\ \iwi{Sanskrit}{\textit{limpáti}} AV+). Second, the extension of the formation to \isi{transitive} verbs denoting cognitive-perceptual and possessional states, which would not a priori be expected from an \isi{anticausative} formant made to \isi{causative} change of state \is{change of state (CoS)} verbs of the sort found here, can be understood to result from either 1) the general reinterpretation and extension of the formant as an all-purpose change of state \is{change of state (CoS)} marker -- e.g., Lith.\ \textit{suprañta, supràsti, supràto}\iw{Lith}{\textit{supràsti, suprañta, supràto}} ‘understand’ (: Go.\ \iwi{Goth}{\textit{fraþjan}} ‘id.’) -- and/or 2) the precocious generalization of the suffix as a regular present-stem formant for (state-oriented) thematic aorists\is{thematic!aorist} (and \isi{stative}/presential perfects),\is{perfect} as normal in later Greek.\footnote{From this starting point the type will then have been fully morphologized as the regular present stem correspondent to thematic aorists\is{thematic!aorist} and then further extended. A similar process will have taken place in Armenian.} 

\section{Indo-Iranian}
\label{sec:IIr}

The type is also unambiguously attested for Indic and indirectly for Iranian -- though here less robustly than for Greek.\footnote{Possible indirect evidence for the type may also survive in the RV and AV in the tendency noted by \citealt{Kulikov2000} for nasal-infix verbs that have relatively well attested athematic and thematic by-forms to show a preference for the thematic variant in \isi{intransitive} usages.} For the \isi{nasal present}s under discussion, see further \citealt{Kulikov2000}, \citealt{Hill2007}, and especially \citealt{Zasada2021} \textit{s.vv}., who provides critical discussion of the relevant forms. The forms can usefully be divided into finite and participial\is{participle} examples. 


\begin{exe}
    \ex Intransitive\is{intransitive} nasal-infix thematic actives in Vedic: Finite Forms
    \begin{xlist}
        \ex { \iwi{Sanskrit}{\textit{tr̥mpáti}} ‘enjoy, take delight in, fill up on (+gen.)’ RV (12x), \iwi{Sanskrit}{\textit{tr̥pṇóti}} ‘id.’ RV (5x), \iwi{Sanskrit}{\textit{tr̥pyati}} AV+: \iwi{Sanskrit}{\textit{átr̥pat}} AVŚ 3.13.6 (: adv. neut. \iwi{Sanskrit}{\textit{tr̥pát}} RV 3x); cf. YAv.\ n.\ sg.\ \iwi{Avestan}{\textit{θra̜fš}°} ‘satisfaction (\textit{vel sim}.)’ Yt.5.26, YAv.\ \iwi{Avestan}{\textit{θra̜fəδra}-} ‘rich in; satisfied', which points to the Indo-Iranian age of the present.}
\ex  \iwi{Sanskrit}{\textit{śrathnās}, \textit{śrathnité}} ‘loosen, make slack’ RV+: \iwi{Sanskrit}{\textit{śrathayanta}} ‘loosen, become slack'.\\

RV 2.24.3ab:\\ \gll
 \textit{tád} \textit{dev{ā́}nāṃ} \textit{devátamāya} \textit{kártuvam} {\textbf{\textit{áśrathnan}} \iw{Sanskrit}{\textit{śrathnās}, \textit{śrathnité}}} \textit{dr̥ḷh{ā́}\textsuperscript{+}} \textit{ávradanta} {\textit{vīḷitā} |} \\
 that.\textsc{nom} god.\textsc{gen.pl} god.\textsc{superl.ins.sg} do.\textsc{ger} loosen.\textsc{3pl.ipf.act} firm.\textsc{nom.pl.n} soften.\textsc{3pl.ipf.mid} hard.\textsc{nom.pl.}\\
 \glt `That had to be done by the foremost god of gods: what was firm became loose, what was hard became pliant.' (\citealt{JamisonBrereton2014})\\

 \vspace{0.3cm} 
 Cf. TS 6.1.9.7 \textit{ánu śr̥nthati} ‘become loose.’ On the geographic localization of the forms of the root with nasal, see \citealt[95\textsuperscript{3}]{Insler1991}.

\ex { \iwi{Sanskrit}{\textit{rṇóti}, \textit{r̥ṇváti}} ‘set in motion, impel (tr., intr. mid.)', \textit{íyarti}, \textit{{ī́}rate} ‘id.’ RV+: \textit{ārta} (intr.) RV.+.}\\

\vspace{0.3cm}
RV 6.2.6ab:\\ 
\gll \textit{tveṣás} \textit{te} \textit{dhūmá} {\textbf{\textit{r̥ṇvati}} \iw{Sanskrit}{\textit{rṇóti}, \textit{r̥ṇváti}}} \textit{diví} \textit{ṣáñ} \textit{chukrá}  {\textit{{ā́}tataḥ} |}\\
turbulent.\textsc{nom.sg} you.\textsc{2sg.gen} smoke.\textsc{nom.sg} impel.\textsc{3sg.prs.act} heaven.\textsc{loc.sg} be.\textsc{prs.ptcp.act.nom.sg} gleaming.\textsc{nom.sg} stretched.out.\textsc{nom.sg}\\
\glt `Your smoke, when it is in heaven, is turbulent in motion, stretched out (there) gleaming, …' (\citealt{JamisonBrereton2014}).\\

\smallskip
{ Possibly further RV 1.144.5cd, 3.11.2c.}

\ex { \iwi{Sanskrit}{\textit{jinóti/te}, \textit{jínvati/te}} ‘revitalize, refresh (tr., intr. mid.)’ RV+}.\\

\smallskip

AVŚ 12.1.46cde:\\ \gll 
\textit{krímir} {\textbf{\textit{jínvat}} \iw{Sanskrit}{\textit{jinóti/te}, \textit{jínvati/te}}} \textit{pr̥thivi}  \textit{yádyad} \textit{éjati} \textit{prāvŕ̥ṣi}  \textit{tán} \textit{naḥ} \textit{sárpan} \textit{mópa} \textit{sr̥pat}\\ 
worm.\textsc{nom.sg} revitalize.\textsc{3sg.ipf.act} Earth.\textsc{voc.sg} whichever move.\textsc{3sg.prs.act} rain.season.\textsc{loc.sg} that.\textsc{nom.sg} us.\textsc{1pl.acc} crawling.\textsc{prs.ptcp.nom.sg.act} \textsc{neg}+towards.\textsc{prvb} crawl.\textsc{3sg.aor.act}\\
\glt `Welcher auflebende Wurm auch immer, o Erde, sich in der Regenzeit regt, der möge, wenn er kriecht, nicht zu uns herankriechen.' \citep[105]{Zasada2021}.\\

\smallskip

{ Possibly also AVŚ 11.4.14c, 16d = AVP 16.22.4b, 6d, AVŚ 12.1.3c. Cf.\ mid.\ 3\textsuperscript{rd} sg.\ \textit{jinvate} (intr.) RV 3.2.11ab. See \citealt[104ff.]{Zasada2021}}

\ex  \iwi{Sanskrit}{\textit{śvíndate}} ‘shine’ Dhātup. (: \textit{áśvitan}, \textit{áśvait}, \textit{śvitāná}{}- RV) (: ORuss.\ \iwi{ORuss}{\textit{svĭ(t)nuti}} ‘grow light', Lith.\ \textit{šviñta, švìsti, švìto}\iw{Lith}{\textit{švìsti, šviñta, švìto}, 1sg. prs. \textit{švintù}} ‘dawn', ON \textit{hvítna}\iw{ON}{\textit{hvítna} (OWN)} ‘goes pale’).\footnote{Given its predominantly active but semantically \isi{intransitive} aorist forms, which points to a verbal root primarily lexified as an \isi{intransitive}, I tentatively include this verb here, despite its likely finite middle inflection, which can easily be explained as analogical \emph{bzw}.\ pleonastic, as outlined below.}
    \end{xlist}
\end{exe} 

\largerpage
\begin{exe}
    \ex  Intransitive\is{intransitive} nasal-infix thematic actives in Vedic: Participial\is{participle} forms
    \begin{xlist}
        \ex { \iwi{Sanskrit}{\textit{siñcáti/te}} ‘pour out (tr., intr. mid.)’ RV+ (: YAv.\ -\textit{hiṇca}-\textit{\textsuperscript{ti} } \iw{Avestan}{-\textit{hiṇca}-} ‘id.’): \iwi{Sanskrit}{\textit{sécate}} ‘pour out (intr.)’ RV 10.96.1 (: OAv.\  PN \iwi{Avestan}{\textit{Haēcat̰}.\textit{aspa}-}), \iwi{Sanskrit}{\textit{sicyáte}} ‘id.’: RV \iwi{Sanskrit}{\textit{ásicat/ta}} RV+}.\\
        
         RV 5.85.6cd:\\
         \gll  \textit{ékaṃ} \textit{yád} \textit{udn{ā́}} \textit{ná} \textit{pr̥ṇánti} \textit{énīr} {\textbf{\textit{ā-siñcántīr}} \iw{Sanskrit}{\textit{siñcáti/te}}} \textit{avánayaḥ} {\textit{samudrám} ||}\\
         single.\textsc{acc.m} which water.\textsc{instr.sg} \textsc{neg} fill.\textsc{3pl.prs.act} mottled.\textsc{nom.pl.f} \textsc{prvb}-pouring.\textsc{prs.ptcp.act.nom.pl.f} stream.\textsc{nom.pl.f} sea.\textsc{acc.sg.m}\\
         \glt `That the mottled streams, pouring out, do not fill the single sea with water.' (\citealt{JamisonBrereton2014}).\\
         \smallskip
         
         { Possible also RV 1.121.6cd (\textit{siñcáñ}). See \citealt[92--95]{Hill2007} with lit.}
\ex { \textit{prá ... r̥ñjánti} \iw{Sanskrit}{\textit{r̥ṇákti}*} ‘stretch out (tr.)', 3\textsuperscript{rd} pl.\ mid.\ \textit{r̥ñjáte} ‘stretch oneself out; move quickly forwards', them.\ \textit{r̥ñjate} ‘id.’: \textit{abhí} ... \iwi{Sanskrit}{\textit{r̥jyate}} RV 1.140.2 ‘id.’}.\\
\smallskip

RV 1.95.7ab:\\
\gll \textit{úd} \textit{yaṃyamīti} \textit{savitéva} \textit{bāhū́}  \textit{ubhé} \textit{sícau} \textit{yatate} \textit{bhīmá}  {\textbf{\textit{r̥ñján}} \iw{Sanskrit}{\textit{r̥ṇákti}*}|}\\ 
up raise.\textsc{3sg.prs.int.act} Savitar.like arm.\textsc{acc.du.m} both.\textsc{acc.du.f} seam.\textsc{acc.du.f} fearsome.\textsc{nom.sg.m} moving.towards.\textsc{3sg.prs.mid} charge.on.\textsc{prs.ptcp.act.nom.sg.m}\\
\glt `Like Savitar, he raises up his arms again. Healing himself along the two seams (of the world?), the fearsome one charging straight on.' (\citealt{JamisonBrereton2014}).\\
\smallskip

RV 1.172.2abc:\\
\gll \textit{āré} \textit{s{ā́}} \textit{vaḥ} {\textit{sudānavo} |} \textit{máruta} {\textbf{\textit{r̥ñjatī́}} \iw{Sanskrit}{\textit{r̥ṇákti}*}} {\textit{śáruḥ} |} \textit{āré} \textit{áśmā} \textit{yám} {\textit{ásyatha} ||}\\
distance.\textsc{loc} \textsc{dem.pron.nom.sg.f} you.\textsc{gen.dat.pl} drop.rich.\textsc{voc.pl} Marut.\textsc{voc.pl} moving.straight.\textsc{prs.ptcp.nom.sg.f} arrow.\textsc{nom.sg.f} distance.\textsc{loc} stone.\textsc{nom.sg} which.\textsc{acc} hurl.\textsc{2pl.prs.act}\\ 
\glt `In the distance be your straight-aiming arrow, you Maruts rich in drops, in the distance the stone that you hurl.' (\citealt{JamisonBrereton2014}).\footnote{Given its \isi{intransitive} semantics, I include this form here; its athematic inflection represents a trivial inner-Vedic innovation.}\\

\newpage
RV 3.31.1:\\ 
\gll \textit{ś{ā́}sad} \textit{váhnir} \textit{duhitúr} \textit{naptíyaṃ} \textit{gād}  \textit{vidv{ā́}m̐} \textit{r̥tásya} \textit{dī́dhitiṃ} {\textit{saparyán} |}  \textit{pit{ā́}} \textit{yátra} \textit{duhitúḥ} \textit{sékam} {\textbf{\textit{r̥ñján}} \iw{Sanskrit}{\textit{r̥ṇákti}*}} \textit{sáṃ} \textit{śagmíyena} \textit{mánasā} {\textit{dadhanvé} ||}\\
instruct.\textsc{prs.ptcp.act.nom.sg.m} conveyor.\textsc{nom.sg.m} daughter.\textsc{gen.sg} (grand)daughter.\textsc{acc.sg} come.\textsc{3sg.aor.inj.act} knowing.\textsc{prf.ptcp.act.nom.sg.m} truth.\textsc{gen.sg} contemplation.\textsc{acc.sg} serving.\textsc{prs.ptcp.act.nom.sg.m} father.\textsc{nom.sg} where daughter.\textsc{gen.sg} outpouring.\textsc{acc.sg} stretching.out.\textsc{prs.ptcp.act.nom.sg} together capable.\textsc{ins.sg} mind.\textsc{ins.sg} run.\textsc{3sg.prf.mid}\\
\glt `Instructing the (grand)daughter [= Āhavanīya?] of his daughter [= Gārhapatya?], the (offering-)conveyor [= Agni] has come -- he the knowing one, ritually serving the visionary power of truth -- to where the father [= priest? Agni?], stretching out straight, has run toward the outpouring of his daughter [=butter offering?] with capable mind.' (\citealt{JamisonBrereton2014}).\\
\smallskip

RV 4.38.7,8:\\ \gll \textit{utá} \textit{syá} \textit{vājī́} \textit{sáhurir} \textit{r̥tā́vā}  \textit{śúśrūṣamāṇas} \textit{tanúvā} {\textit{samaryé} |} \textit{túraṃ} \textit{yatī́ṣu} \textit{turáyann} \textit{r̥jipyó} \textit{ádhi} \textit{bhruvóḥ} \textit{kirate} \textit{reṇúm}  {\textbf{\textit{r̥ñján}} \iw{Sanskrit}{\textit{r̥ṇákti}*}{||}}\\
and \textsc{dem.pron.nom.sg.m} prizewinner.\textsc{nom.sg.m} victorious.\textsc{nom.sg.m} truthful.\textsc{nom.sg.m} fame.seeking.\textsc{des.pctp.mid.nom.sg.m} body.\textsc{instr} clash.\textsc{loc} headlong.\textsc{adv} go.\textsc{prs.ptcp.act.loc.pl.f} hasten.\textsc{prs.act.pctp.nom.sg.m} hurrying.\textsc{nom.sg.m} towards eyebrows.\textsc{loc.du.f} scatter.\textsc{prs.mid.3sg} dust.\textsc{acc} stretch.out.\textsc{prs.ptcp.act.nom.sg.m}\\

\gll \textit{utá} \textit{sma} \textit{asya} \textit{tanyatór} \textit{iva} \textit{dyór}  \textit{r̥ghāyató} \textit{abhiyújo} {\textit{bhayante} |} \textit{yad{ā́}} \textit{sahásram} \textit{abhí} \textit{ṣīm} \textit{áyodhīd} \textit{durvártuḥ} \textit{smā} \textit{bhavati} \textit{bhīmá} {\textbf{\textit{r̥ñján}} \iw{Sanskrit}{\textit{r̥ṇákti}*}||}\\
and \textsc{emph} \textsc{dem.pron.gen.sg.m} thunder.\textsc{abl} like heaven.\textsc{gen} tremble.\textsc{prs.pctp.act.gen.sg} charge.\textsc{abl} fear.\textsc{prs.mid.3pl} when thousand to \textsc{pers.pron.acc.sg} battle.\textsc{aor.act.3sg} difficult.to.obstruct.\textsc{nom.sg.m} \textsc{emph} become.\textsc{prs.act.3sg} fearsome.\textsc{nom.sg.m} stretch.out.\textsc{prs.ptcp.act.nom.sg.m}\\

\glt `And this prizewinner, victorious, truthful, himself seeking fame with his own body in the clash, hastening headlong toward them (fem.\  [=mares?]) as they go hastily, straight-flying, scatters dust up to the eyebrows as he stretches out straight. And they take fear at his charge as he shows his mettle, as if at the thundering of heaven. When a thousand have battled him, the fearsome one becomes difficult to obstruct, as he stretches out straight.' (\citealt{JamisonBrereton2014}).
\ex { \iwi{Sanskrit}{\textit{r̥ṇádhat}} ‘accomplish’ RV AV, \iwi{Sanskrit}{\textit{r̥dhnóti}} ‘id.’ RV 1x, AV 1x: \iwi{Sanskrit}{\textit{r̥dhyate}/\textit{r̥dhyáte}} ‘succeed’ RV+}.\\
\smallskip

 RV 1.173.11:\\ \gll \textit{yajñó} \textit{hí} \textit{ṣma} \textit{índaraṃ}\textsuperscript{+} \textit{káś cid} {\textbf{\textit{r̥ndháñ}} \iw{Sanskrit}{\textit{r̥ṇádhat}}} \textit{juhurāṇáś} \textit{cin} \textit{mánasā} {\textit{pariyán} |} \textit{tīrthé} \textit{ná} \textit{áchā} \textit{tātr̥ṣāṇám} \textit{óko}  \textit{dīrghó} \textit{ná} \textit{sidhrám} \textit{ā́} \textit{kr̥ṇoti} {\textit{ádhvā} ||}\\
 sacrifice.\textsc{nom.sg.m} for \textsc{emph} Indra.\textsc{acc} whichever.\textsc{nom.sg.m} suceeding.\textsc{prs.ptcp.act.nom.sg.m} swerving.\textsc{aor.ptcp.mid.nom.sg.m} \textsc{emph} mind.\textsc{ins.sg} go.around.\textsc{prs.ptcp.act.nom.sg.m}  ford.\textsc{loc.sg} as to thirst.\textsc{prf.ptcp.mid.acc.sg.m} home.\textsc{acc.sg.n} long.\textsc{nom.sg.m} as goal.\textsc{acc.sg} \textsc{prvb} do.\textsc{3sg.prs.act} road.\textsc{nom.sg.m} \\
\glt   `For any sacrifice that reaches fulfillment, even though it swerves along, meandering in mind, brings Indra to the house, as if bringing a thirsting man to a ford -- as a long road brings home a man who reaches his goal.' (\citealt{JamisonBrereton2014}).\\

\ex  \iwi{Sanskrit}{\textit{min{ā́}ti/te}} ‘damage, harm; exchange, change (tr., intr. mid.)’ RV+ (: \iwi{Sanskrit}{\textit{m{ī́}yate}} ‘disappear, pass away’ RV).\\
\smallskip

RV 5.2.1cd:\\ \gll \textit{ánīkam} \textit{asya} \textit{ná} {\textbf{\textit{mináj}} \iw{Sanskrit}{\textit{min{ā́}ti/te}}} \textit{jánāsaḥ}  \textit{puráḥ} \textit{paśyanti} \textit{níhitam} {\textit{arataú} |}\\
face.\textsc{nom.sg.n} his \textsc{neg} changing.\textsc{prs.ptcp.act.nom.sg.n} people.\textsc{nom.pl} front see.\textsc{3pl.prs.act} down.set.\textsc{acc.sg.n} spokes.\textsc{loc}\\
\glt `His face is not one that changes (its face): the peoples see it in front, set down in the circle of spokes.' (\citealt{JamisonBrereton2014}).\\

\smallskip
Cf.\ intr.\ them.\ mid.:\\\smallskip

RV 1.79.2a:\\ \gll 
\textit{{ā́}} \textit{te} \textit{suparṇ{ā́}} {\textbf{\textit{aminantam̐}} \iw{Sanskrit}{\textit{min{ā́}ti/te}}} \textit{évaiḥ}\\
\textsc{prvb} your.\textsc{gen.sg} fine.feathered.\textsc{nom.pl.m} change.\textsc{3pl.ipf.mid} way.\textsc{ins.pl}\\
\glt `Your fine-feathered (lightning flashes) zigzagged along their ways.' (\citealt{JamisonBrereton2014}).\\

\ex  \iwi{Sanskrit}{\textit{muñcáti/te}} ‘release, set free (tr., intr. mid.)’ RV+: \iwi{Sanskrit}{\textit{múcyate/mucyáte}} ‘get free’ RV+ (: \iwi{Sanskrit}{\textit{ámucat}} RV+, {\textit{ámugdhvam}} RV+; OCS \textit{promŭknǫ, -mŭknǫti sę},\iw{OCS}{\textit{promŭknǫti sę, -mŭknǫ, -mŭkŭ}} Lith.\ \textit{muñka, mùkti}\iw{Lith}{\textit{mùkti, muñka, mùko}}).\footnote{Note further OP <a-m\textsuperscript{u}{}-u-θ> /amunθa/ \iw{OP}{\textit{amunθa}} D I 01, ii 2, 71, iii 71 ‘fled, escaped', which as generously pointed out by a reviewer is most easily taken as an \isi{intransitive} thematic imperfect.}\\\smallskip

AVŚ 8.7.10ab:\\ \gll 
{\textbf{\textit{unmuñcántīr}} \iw{Sanskrit}{\textit{muñcáti/te}}} \textit{vivaruṇ{ā́}} \textit{ugr{ā́}} \textit{y{ā́}} {\textit{viṣadū́ṣanīḥ} |}\\
\textsc{prvb}.get.free.\textsc{prs.ptcp.nom.pl.f} off.Varuṇa.\textsc{nom.pl.f} mighty.\textsc{nom.pl.f} who.\textsc{nom.pl.f} destroying.poison.\textsc{nom.pl.f}\\
\glt  `Sich befreiend, frei von Varuṇa, mächtig, die Gift zerstörend sind... [= Pflanzen].' \citep[214]{Zasada2021}.\\\smallskip
 
 {  Cf. 8.7.4 (intr.) \iw{Sanskrit}{\textit{prastr̥ṇatī}-} \textit{prastr̥ṇatīs},  \iw{Sanskrit}{\textit{pratanvatī}-} \textit{pratanvatīs}. For full discussion, see \citealt[214]{Zasada2021}.}
 
\ex { \iwi{Sanskrit}{\textit{ásinvant}-} ‘insatiable’ RV (5x; \textit{asinvá}- ‘id.’ RV) if as likely to tr.\ \iwi{Sanskrit}{*\textit{sinóti}} ‘satiate’ (: TA \iwi{TA}{\textit{sim̐señc}} ‘satiate', mid.\ \textit{sim̐santär} ‘become satiated', TB \iwi{TB}{\textit{sinastar}} ‘id.’)}.\\\smallskip

RV 10.79.1d:\\ \gll 
{\textbf{\textit{ásinvatī}} \iw{Sanskrit}{\textit{ásinvant}-}} \textit{bápsatī} \textit{bhū́ri} \textit{attaḥ}\\
insatiable.\textsc{nom.du.f} gnawing.\textsc{prs.ptcp.nom.du.f} much.\textsc{acc.sg} eat.\textsc{3du.prs.act}\\
\glt `Insatiable, gnawing, they eat amply.' (\citealt{JamisonBrereton2014}).
\end{xlist}
\end{exe}


There are two interesting observations that can be made about this material.\footnote{It is also important to note that this type regularly correlates with the thematic aorist\is{thematic!aorist} in Indic and Iranian -- cf., e.g., \iwi{Sanskrit}{\textit{tr̥mpáti}} ‘enjoy, take delight in, fill up on (+gen.)’ RV: \textit{átr̥pat} AVŚ 3.13.6 (: adv. neut. \textit{tr̥pát} RV) or \iwi{Sanskrit}{\textit{siñcáti/te}} ‘pour out (tr., intr. mid.)’ RV+ (: YAv.\ -\textit{hiṇca}-\textit{\textsuperscript{ti}} \iw{Avestan}{-\textit{hiṇca}-} ‘id.’): RV \iwi{Sanskrit}{\textit{ásicat/ta}} RV+ -- as regularly in Greek (and Armenian) and systematically in OCS. This is obviously an inherited feature of a subset of the verbs in this class.} First, there is a fundamental difference in the distributional attestation of this type in Greek and Indo-Iranian. In Greek, nearly all the forms we find are independent lexical items that synchronically lack the corresponding \isi{transitive} nasal-infix form from which they have been derived. In Indo-Iranian, while there are such cases, we mostly find that the forms are independent, one-off active thematic \isi{intransitive} uses of an otherwise well-attested nasal stem that normally signals causativity\is{causative} alternations with active vs. middle voice. This distribution suggests that the derivational type survived as a living process into Indo-Iranian and then early Indic at least in some verbs and/or as a dialectal phenomenon. What we have in our texts should thus likely be considered a mix of archaic and/or dialect forms and poetic innovations based on them.\footnote{The poetic character of the type can be seen clearly in the formulaic-template-type usage of, e.g., \textit{pāda}-final \textit{r̥ñján}. \iw{Sanskrit}{\textit{r̥ṇákti}*}}


\begin{table}
\caption{Derivational \& chronological relationship between athematic and thematic nasal-infix stems}
    \begin{tabularx}{\textwidth}{p{6.3cm}cX}
\lsptoprule
{\bfseries Nasal-Infix Causative}\is{causative} & & {\bfseries Nasal-Infix Anticausative}\is{anticausative}\\
\midrule
{ PIIr.\textsubscript{1} \iwi{PIIr}{*\textit{sinákti, sinktái}} ‘pour out (tr., intr.)’ } &            → &  *\textit{sinčáti} ‘pour out (intr.)’  \\
& &                 \multicolumn{1}{c}{$\Downarrow$} \\
\multicolumn{1}{c}{$\Downarrow$}  & & \mbox{}{ \iwi{PIIr}{*\textit{sinčáti}, \textit{sinčátai}} ‘id.’}\\
& &                  \multicolumn{1}{c}{$\Downarrow$}   \\
 PPIr.\textsubscript{3} \iwi{PIIr}{*\textit{sinčáti}, \textit{sinčátai}} ‘pour out (tr., intr.)’ {            (> Ved.\ \iwi{Sanskrit}{\textit{siñcáti/te}}, YAv.\ \textit{hiṇca}-\textit{\textsuperscript{ti}})} \iw{Avestan}{-\textit{hiṇca}-} & $\Leftarrow$ & *\textit{sinčátai} ‘id.’ (residual \iwi{PIIr}{*\textit{sinčánt}-} > Ved.\ \iwi{Sanskrit}{\textit{siñcánt}-} ‘id.’)\\
\lspbottomrule
    \end{tabularx}
    \label{tab:deriv}
\end{table}


\largerpage
The second thing to note is that the preservation of this type into Indo-Iranian provides a straightforward explanation of another remarkable feature of some of the verbs that show this usage -- the fact that already in the RV and in some cases in Avestan they are fully \is{thematization}thematized, without any trace of their original athematic inflection -- viz.\ \iwi{Sanskrit}{\textit{siñcáti/te}} RV+: YAv.\ -\textit{hiṇca}-\textit{\textsuperscript{ti}}, \iw{Avestan}{-\textit{hiṇca}-} \iwi{Sanskrit}{\textit{muñcáti/te}} RV+ (: OCS \textit{promŭknǫ, -mŭknǫti sę},\iw{OCS}{\textit{promŭknǫti sę, -mŭknǫ, -mŭkŭ}} Lith.\ \textit{muñka, mùkti}\iw{Lith}{\textit{mùkti, muñka, mùko}}). For the attestation of these forms and others like it, see \citealt{Hill2007}. The highly precocious \isi{thematization} we find here and in similar cases -- so, e.g., \iwi{Sanskrit}{\textit{kr̥ntáti/te}} RV+, YAv.\ \textit{kərəṇta-\textsuperscript{ti/te}} \iw{Avestan}{\textit{kərəṇta}-} (: OHG \mbox{\iwi{OHG}{\textit{scrintan}})} and other stems with obstruent final roots liable to the \isi{causative} alteration -- makes perfect sense under the assumption that in some cases the thematic \isi{intransitive} active was subject to an early pleonastic medialization, which then led to the \isi{backformation} of a thematic \isi{transitive} active and the displacement of the original athematic stem. The full derivational and chronological relationship between these forms is given in \tabref{tab:deriv}.



Both analogical developments are well paralleled in Indic and Iranian -- cf., e.g., Ved.\ \iwi{Sanskrit}{\textit{édhate}} ‘flourish, increase’ (: {αἴθω} \iw{AG}{αἴθω, -oμαι} ‘burn (intr.)’ Pi.S., \iwi{AG}{αἴθoυσα} ‘portico’ Hom.+;  αἴθoμαι Hom.+, ClArm. \iwi{Armenian}{\textit{ayrem}} ‘burn’), for “pleonastic” medialization, and Ved.\ \iw{Sanskrit}{\textit{śráyati/te}} \textit{śráyati} ‘lean on (tr.)’ (: \iw{Avestan}{-\textit{sirinao}-} YAv.\ -\textit{sirinao}-\textit{\textsuperscript{ti}}; Ved.\ \textit{śráyate}), for the \isi{backformation} of a \isi{transitive} active from an \isi{intransitive} middle stem.\footnote{See further footnote \ref{fn:itract}.}

\section{Conclusion}

Indo-Iranian and Ancient Greek\footnote{Further Classical Armenian. See footnote \ref{fn:ClassArm}.}  thus provide a very close parallel to the type outlined by Gorbachov for the \isi{Northern Indo-European} languages -- one that matches it morphologically, semantically, and derivationally. Based on the close agreement of these languages, it is possible to confirm Gorbachov’s analysis of the \isi{Northern Indo-European} material and to reconstruct the formation for late Nuclear Proto-Indo-European as an important deverbal\is{deverbal/deverbative!verb} and deadjectival\is{deadjectival!verb}  \isi{intransitive} change of state\is{change of state (CoS)} type.\footnote{The type is also attested for Italo-Celtic as \citet[217]{Gorbachov2007}, following others, rightly notes -- cf., e.g., Lat.\ \iwi{Latin}{°\textit{cumbō}}, -\textit{ere} Pl.+, \iwi{Latin}{\textit{ning(u)it}}  Lucr.+, and possibly  \iwi{Latin}{°\textit{hendō}}, -\textit{ere} Pl.+. A task for the future is to identify further instances in these branches.} 

\section*{Abbreviations}
\begin{tabularx}{.45\textwidth}{lQ}
A.R. & Apollonius Rhodius \\
Alb. & Albanian \\
Ar. & Aristophanes \\
Arc. & Arcadian \\
Archil. & Archilochus \\
Arist. & Aristotle \\
Att. Ion. & Attic Ionic \\
AV & Atharvaveda \\
AVP & Atharvaveda Paippalada \\
AVS & Atharvaveda Saunakiya \\
Call. & Callimachus \\
ClArm. & Classical Armenian \\
Dhātup. & Dhātupātha \\
Dial. & Dialects \\
\end{tabularx}
\begin{tabularx}{.45\textwidth}{lQ}
E. & Euripides \\
Eust. & Eustathius \\
Gal. & Galen \\
Gk. & Greek \\
Go. & Gothic \\
Hdt. & Herodotus \\
Hes. & Hesiod \\
Hes. Sc. & Hesiod Scutum \\
hHom. & Homeric Hymn \\
Hitt. & Hittite \\
Hom. & Homer \\
Hp. & Hippocrates \\
Hsch. & Hesychius\\
Intr. & Intransitive \\
KS & Katha Samhita \\
Lat. & Latin \\
\end{tabularx}

\begin{tabularx}{.45\textwidth}{lQ}
Lesb. & Lesbian \\
Lith.\ & Lithuanian \\
Lucr. & Lucretius \\
MW & Middle Welsh \\
Myc. & Mycenaean \\
OAv. & Old Avestan \\
OCS & Old Church Slavonic \\
OE & Old English \\
OHG & Old High German \\
OIr. & Old Irish \\
ON & Old Norse \\
OP & Old Persian\\
OPr.\ & Old Prussian \\
ORuss. & Old Russian \\
PGmc. & Proto-Germanic \\
\end{tabularx}
\begin{tabularx}{.45\textwidth}{lQ}
Pi. & Pindar \\
Pl. & Plautus \\
PN & Proper name \\
RV & Rigveda \\
S. & Sophocles \\
Sa. & Sappho \\
TA & Tocharian A \\
TB & Tocharian B \\
Th. & Thucydides \\
Theocr. & Theocritus \\
Tr. & Transitive \\
TS & Taittiriya Samhita \\
Ved. & Vedic Sanskrit \\
YAv. & Young Avestan \\
YV & Yajurveda \\
\end{tabularx}



{\sloppy\printbibliography[heading=subbibliography,notkeyword=this]}
\end{document}

